\documentclass[11pt,a4paper]{article}
\usepackage[margin=28mm]{geometry}
\usepackage{amsmath,amssymb,amsthm,mathtools,booktabs,longtable,array,enumitem}
\usepackage[T1]{fontenc}
\usepackage{lmodern,microtype}
\usepackage[hidelinks]{hyperref}
\usepackage{xurl}
\allowdisplaybreaks[2]
\setlength{\emergencystretch}{2em}
\newtheorem{theorem}{Theorem}[section]
\newtheorem{proposition}[theorem]{Proposition}
\newtheorem{lemma}[theorem]{Lemma}
\newtheorem{corollary}[theorem]{Corollary}
\theoremstyle{definition}
\newtheorem{definition}[theorem]{Definition}
\newtheorem{remark}[theorem]{Remark}
\newtheorem{sourceinput}[theorem]{Source input}
\newcommand{\C}{\mathbb C}
\newcommand{\Q}{\mathbb Q}
\newcommand{\Z}{\mathbb Z}
\newcommand{\F}{\mathbb F}
\newcommand{\N}{\mathbb N}
\newcommand{\HH}[2]{\mathsf H_{#1}(#2)}
\newcommand{\vnorm}[1]{\left\lVert#1\right\rVert}
\newcommand{\abs}[1]{\left|#1\right|}
\newcommand{\ip}[2]{\left\langle#1,#2\right\rangle}
\newcommand{\Aut}{\operatorname{Aut}}
\newcommand{\Span}{\operatorname{span}}
\newcommand{\im}{\operatorname{im}}
\newcommand{\length}{\operatorname{length}}
\newcommand{\diag}{\operatorname{diag}}
\newcommand{\volsq}{\operatorname{vol}^{2}}
\newcommand{\src}[2]{\par\smallskip\begingroup\footnotesize\noindent Source: \path{#1}.\par\noindent\path{#2}.\par\endgroup}
\title{An exposition of the compact RH and GRH development\\[4pt]
\large Compatible twists, denominators, and the volume inequality}
\author{A working exposition prepared for Christian Kudeba\\
with reductions from Principia contributors}
\date{27 November 2026}
\begin{document}
\maketitle
\begin{abstract}
We develop a prose account of the compact formal development at commit
\texttt{72fead102f663d5c7a407964a13dc0cf9d977207} of
\texttt{lsernau/grh-lean}, directory \texttt{grh-compact}.
This revision starts with the arithmetic inputs and Theorem C of its three-theorem reduction.
It gives full proofs of the degree-independent polynomial estimate used to choose compatible
powers of a twist, the determinant-specialization argument for finite relation modules,
the linear algebra that converts scalar operators into a denominator bound, the finite mirror estimate, the global reciprocity calculation, bounded pencils of tied primes,
the simultaneous compatible-twist selection, the numerical budget comparisons, and the
height and Gram-determinant calculus used at the end of Theorem C.
The local balancing and special-value verification for automorphic rows, the Ribet and patching operators,
and the geometric comparison forms are explicitly identified as inputs still to be
expanded here. Thus this revision is a partial exposition, not a standalone prose proof
of RH. The formal development contains proofs of these inputs; an unexpanded input in
this text does not mean a missing lemma in that development.
\end{abstract}
\tableofcontents

\section{The edition and the argument}
The proof under discussion is attributed to Levent Alp\"oge and Claude.
The compact edition is due to Luke Sernau, colleagues at Anthropic, and Claude.
We use Bartosz Naskr\k{e}cki's working reading as a guide, at commit
\texttt{82252d4e71d9f4db30a38e370dba78283ab3a921} of
\texttt{nasqret/riemann-rev}. Its text is available under 0BSD.
Relevant portions of the notes \emph{Short Selmer groups}, \emph{Fit twists}, \emph{Heights},
\emph{Chain volume}, and \emph{Lift volume height} were compared with the corresponding
pinned Lean files for this revision. Their reference to an input or an outline is
not treated as a complete proof of that input. The proofs below are rewritten in one
notation; source pointers at the end identify their formal counterparts.

The compact snapshot on Principia records a compilation of its 794 modules and audits
of its five concluding declarations against the three standard axioms
\texttt{propext}, \texttt{Classical.choice}, and \texttt{Quot.sound}.
Its receipt does not record a full Comparator pass. These are statements about the
formal artifact and its recorded checks; they do not certify the correctness or
completeness of the present prose.

The three structural statements have the following roles. Theorem A supplies
many short theta lifts, with a bound for the index needed to make them integral.
Theorem B says that a logarithmically bounded number of their translates almost spans
the row space. Theorem C bounds below the covolume of any such almost-spanning,
small-index family. Comparing the two estimates bounds a weighted sum over non-split
primes. A good-window argument converts the bound to a prime-error estimate, and an
analytic nonvanishing criterion gives the five concluding statements about RH and GRH.

The present revision is organized around the hard interface between A and C:
\[
 \begin{gathered}
 \text{short twisted Selmer groups}\ \longrightarrow\
       \text{projection denominators}\\
 \longrightarrow\ \text{height of volume ratios}
       \ \longrightarrow\ \text{Theorem C}.
 \end{gathered}
\]
A compatible twist must work simultaneously for a finite set of Galois rows and
for each of its first \(m\) powers. It is chosen separately for each prime and each
identification of the complex numbers with an algebraic closure of a local field.
One global complex character working at every prime is neither asserted nor needed.

\paragraph{What carries over from the earlier manuscript.}
The existing paper supplies a publication record, source history, and material for reuse.
The analytic strategy, explanations of determinant comparisons, and much of the local
terminology can be reused after checking their definitions. Earlier fourth-power
windows, auxiliary constants, and conditional character estimates are not silently
inserted into this compact route. The earlier 668-page edition remains available in the
\href{https://github.com/ChristianKudeba/alpoge-rh-verification/releases/tag/short-route-6be17bf2cd8bf8b790193532ca4d7e4616ba5b45}{source archive}.
Its page count measures written material, not verified coverage of the compact proof.

\paragraph{Conventions.}
All logarithms of positive real numbers are natural. For a prime \(\ell\),
\(\overline{\Q}_\ell\) carries the absolute value \(\abs{\ell}_\ell=\ell^{-1}\);
write \(v_\ell(z)=-\log\abs z_\ell/\log\ell\), with \(v_\ell(0)=+\infty\).
The valuation ring of a valued field \(K\) is \(\mathcal O_K\).
Inner products are conjugate-linear in the first variable and linear in the second.
For a linear map \(T\) from a finite-dimensional Hermitian space put
\(\volsq(T)=\det(T^*T)\). All products or sums over primes below a bound are finite.
A \emph{source input} is a mathematical assertion whose proof has not yet been
fully written in this manuscript.


\section{A uniform estimate at roots of unity}\label{sec:torsion}
This is the polynomial part of the compatible-twist construction.
The order of the root of unity is chosen before the polynomials are given.
There is no bound on their degrees.

\subsection{Gauss norms}
For a polynomial \(F(X)=\sum_i a_iX^i\) over a non-archimedean valued field and
a real number \(0<\rho\le1\), put
\[
 \vnorm{F}_\rho=\max_i \abs{a_i}\rho^i,\qquad \vnorm0_\rho=0.
\]
This norm is non-archimedean, multiplicative, and nondecreasing as a function of
\(\rho\). Here is a proof of the only less immediate assertion, multiplicativity.
Choose the largest indices \(i,j\) at which the respective maxima for two nonzero
polynomials \(F,G\) are attained. In the coefficient of \(X^{i+j}\), the summand
\(a_i b_j\), after multiplication by \(\rho^{i+j}\), has absolute value
\(\vnorm F_\rho\vnorm G_\rho\). Every other summand is strictly smaller:
if its first index exceeds \(i\), its first factor is strictly smaller than the
first maximum; otherwise its second index exceeds \(j\). The ultrametric
inequality, with one strictly largest summand, gives equality in this coefficient.
The reverse bound for every coefficient follows by the ultrametric inequality.
Also, if \(\abs z\le\rho\), then \(\abs{F(z)}\le\vnorm F_\rho\).

If \(\eta\) is a primitive root of unity of order \(\ell^{j+1}\), then
\begin{equation}\label{eq:torsion-radius}
 \abs{\eta-1}=\abs\ell^{1/(\ell^j(\ell-1))}.
\end{equation}
Indeed all \(\ell\)-power roots reduce to \(1\) in residue characteristic \(\ell\).
For \(u\) prime to \(\ell\), the quotient
\((\eta^u-1)/(\eta-1)\) reduces to \(u\), and is a unit.
All primitive roots therefore have the same distance from \(1\).
Their distances multiply to \(\abs{\Phi_{\ell^{j+1}}(1)}=\abs\ell\), and
there are \(\ell^j(\ell-1)\) of them. Roots of lower order are no farther from \(1\).
The cyclotomic identity here follows directly from
\(\Phi_{\ell^{j+1}}(X)=\sum_{a=0}^{\ell-1}X^{a\ell^j}\).

\begin{theorem}[degree-independent simultaneous evaluation]\label{thm:torsion}
Let \(\ell\) be prime, \(0<r<1\), and \(\abs\ell<r^{\ell-1}\).
Fix a finite set \(I\) and positive real numbers \(y_i\), \(i\in I\).
There is \(N_0\), depending only on these data, such that for every \(N\ge N_0\)
the following holds in every valued extension \(K\) of \(\Q_\ell\) containing
the \(\ell^N\)-th roots of unity. If
\(F_i\in\mathcal O_K[X]\) and \(\vnorm{F_i}_r\ge y_i\) for all \(i\), there is
\(\zeta^{\ell^N}=1\) with
\[
 \abs{F_i(\zeta-1)}\ge r y_i\qquad(i\in I).
\]
\end{theorem}
\begin{proof}
The assertion is immediate when \(I\) is empty. Since integral polynomials have
Gauss norm at most \(1\) at a radius at most \(1\), the other hypotheses are
impossible if any \(y_i>1\). We may assume \(0<y_i\le1\).
Put \(Y=\prod_i y_i>0\), and choose an integer \(a\ge1\) with \(r^a<\abs\ell\).
Choose \(N_0\) so that, for every \(N\ge N_0\),
\begin{equation}\label{eq:harmonic-threshold}
 \frac{(\ell-1)(-\log r)}a\sum_{j=0}^{N-1}\frac1{j+1}>-\log Y .
\end{equation}
Such a threshold exists because the harmonic series diverges. For completeness,
its terms with \(2^s\le j+1<2^{s+1}\) sum to at least \(1/2\);
arbitrarily many disjoint blocks give arbitrarily large sums.

Suppose a family \(F_i\) is a counterexample for \(N\ge N_0\). Write
\[
 \Pi=\prod_iF_i,\qquad
 \rho_t=r^{1/\ell^t},\qquad G_t=\vnorm{\Pi}_{\rho_t}.
\]
Multiplicativity and monotonicity give
\(1\ge G_t\ge\prod_i\vnorm{F_i}_r\ge Y>0\).

Fix \(0\le j<N\), put \(n=\ell^{j+1}\), \(p=a(j+1)\), and
\(P(X)=(1+X)^n-1\). Its distinct roots \(\alpha\) are the numbers \(\eta-1\)
with \(\eta^n=1\). By \eqref{eq:torsion-radius} and the hypothesis on \(r\),
every such root has \(\abs\alpha<\rho_j\le\rho_t\) for \(t=j,j+1\).
For each root the assumed failure supplies an \(i\) with
\(\abs{F_i(\alpha)}<r y_i\). Consequently, for \(t=j,j+1\),
\begin{equation}\label{eq:product-small-at-root}
 \abs{\Pi(\alpha)}<r y_i\prod_{k\ne i}\vnorm{F_k}_{\rho_t}
 \le r G_t .
\end{equation}

Divide \(\Pi^p\) by the monic polynomial \(P\):
\(\Pi^p=PS+R\), where \(\deg R<n\).
Interpolation at the roots of \(P\) gives
\[
 R(X)=\sum_{P(\alpha)=0}\Pi(\alpha)^p
               \frac{P(X)}{P'(\alpha)(X-\alpha)} .
\]
To verify this formula, evaluate it at each root: exactly the corresponding term
survives, with value \(\Pi(\alpha)^p\); two polynomials of degree less than \(n\)
with the same \(n\) values are equal. Since
\(\abs{P'(\alpha)}=\abs{n(1+\alpha)^{n-1}}=\abs\ell^{j+1}\), we have
\[
 \vnorm R_{\rho_t}
 < r^pG_t^p\abs\ell^{-(j+1)}\rho_t^{\,n-1}
 \le \left(\frac{r^a}{\abs\ell}\right)^{j+1}G_t^p
 <G_t^p.
\]
Indeed \(P/(X-\alpha)\) is a product of \(n-1\) linear factors whose roots
have absolute value at most \(\rho_t\), hence has Gauss norm
\(\rho_t^{n-1}\). The strict estimate for each summand remains strict for their
finite maximum.

As \(\vnorm{\Pi^p}_{\rho_t}=G_t^p\) and the remainder is strictly smaller,
the ultrametric inequality forces
\[
 \vnorm S_{\rho_t}=G_t^p/\rho_t^n .
\]
Monotonicity of the Gauss norm of this same quotient \(S\) at the two radii gives
\[
 \log G_{j+1}-\log G_j
 \ge \frac n p(\log\rho_{j+1}-\log\rho_j)
 =\frac{(\ell-1)(-\log r)}{a(j+1)}.
\]
Sum over \(0\le j<N\). The left side is
\(\log G_N-\log G_0\le-\log Y\), contradicting
\eqref{eq:harmonic-threshold}. All choices were independent of the degrees and
coefficients of the \(F_i\).
\end{proof}
\src{New/Selmer/RootsOfUnityPoints.lean}{exists_root_of_unity_eval_large}

This proof makes explicit the Gauss-norm computation outlined in
Naskr\k{e}cki's \emph{Short Selmer groups}, Lemma on one variable.
It uses a Lagrange remainder in place of the equivalent differentiated product
identity in the formal proof. It proves the same degree-independent estimate;
it is not a claim that the formal file has been shortened.

\subsection{Several variables and compatible powers}
\begin{theorem}[one tuple for all powers]\label{thm:many-torsion}
Fix a prime \(\ell\), positive integers \(a,m\), nonnegative integers \(q,L\),
a finite set \(I\), nonnegative real numbers \(B_i\), and \(\delta>0\).
There is \(N_0\) depending only on these data with the following property.
Let \(N\ge N_0\), and let \(K\) be a valued extension of \(\Q_\ell\) containing
the \(\ell^N\)-th roots of unity. Let \(S,T\) be finite sets with \(\#S\le q\).
Suppose
\[
 D_t\in\mathcal O_K[X]\text{ is monic},\quad
 \sum_{t\in T}\deg D_t\le L,
\]
and, for \(i\in I\), \(1\le d\le m\), suppose
\(F_{i,d}\in\mathcal O_K[X_s:s\in S]\) and
\(u_{i,d}\in\mathcal O_K^S\) satisfy
\[
 \abs{u_{i,d,s}-1}^{a}\le\abs\ell\quad(s\in S),
 \qquad v_\ell(F_{i,d}(u_{i,d}))\le B_i .
\]
Then some \(\zeta\in\mathcal O_K^S\) satisfies
\[
 \zeta_s^{\ell^N}=1,\qquad
 v_\ell(F_{i,d}((\zeta_s^d)_s))\le B_i+\delta,\qquad
 v_\ell(D_t(\zeta_s))\le\delta
\]
for every indicated index.
\end{theorem}
\begin{proof}
Choose \(0<r<1\) sufficiently close to \(1\) that
\begin{equation}\label{eq:many-r}
 \max(\abs\ell,\ell^{-\delta})
       <r^{\,q+L+\ell+ma+1}<1 .
\end{equation}
In particular \(\abs\ell<r^{\ell-1}\), \(\abs\ell<r^{ad}\) for \(1\le d\le m\),
and \(r^q,r^{L+1}>\ell^{-\delta}\).
Apply Theorem~\ref{thm:torsion} in advance, for each \(0\le h\le q\), to the
finite list of lower bounds
\[
 r^L,\qquad \ell^{-B_i}r^h\quad(i\in I,\ 1\le d\le m).
\]
Take the maximum of these finitely many thresholds for \(N_0\).

We replace the coordinates of each \(u_{i,d}\) one at a time.
After replacing the coordinates of a subset \(A\subseteq S\), require that
the value of each \(F_{i,d}\) at the mixed tuple is at least
\(\ell^{-B_i}r^{\#A}\) in absolute value. The assertion holds for \(A=\varnothing\).

Choose \(s\notin A\). For every \(i,d\), keep the other mixed coordinates fixed
and substitute \((1+X)^d\) for coordinate \(s\), obtaining \(H_{i,d}(X)\).
Its Gauss norm at \(r\) is at least \(\ell^{-B_i}r^{\#A}\).
To see this, the monic polynomial
\[
 ((1+X)^d-u_{i,d,s})^a
\]
has degree \(ad>0\), and its constant coefficient has absolute value at most
\(\abs\ell<r^{ad}\). In an algebraic closure its constant coefficient is the
product, up to sign, of its \(ad\) roots. Thus some root \(w\) has
\(\abs w<r\). Evaluation at \(w\) sets the coordinate back to \(u_{i,d,s}\)
and gives the old mixed value; the Gauss bound on evaluation proves the assertion.

Add the polynomial \(H_0(X)=\prod_{t\in T}D_t(1+X)\) to this finite list.
It is monic of degree at most \(L\), so \(\vnorm{H_0}_r\ge r^L\).
Theorem~\ref{thm:torsion}, with \(h=\#A\), supplies a root
\(\zeta_s^{\ell^N}=1\) such that
\[
 \abs{H_{i,d}(\zeta_s-1)}
       \ge\ell^{-B_i}r^{\#A+1},
 \qquad
 \abs{\prod_tD_t(\zeta_s)}\ge r^{L+1}.
\]
Every factor \(D_t(\zeta_s)\) has absolute value at most \(1\).
It follows that each factor has absolute value at least \(r^{L+1}>
\ell^{-\delta}\). Keep \(\zeta_s\), and continue the induction.

After at most \(q\) steps all coordinates have been replaced. Each value of
\(F_{i,d}\) has absolute value at least
\(\ell^{-B_i}r^{\#S}\ge\ell^{-B_i}r^q>\ell^{-(B_i+\delta)}\).
Taking valuations proves the claims. If \(S\) is empty the assertion follows
directly from the hypotheses.
\end{proof}
\src{New/Selmer/RootsOfUnityPoints.lean}{exists_roots_of_unity_ellVal_le}

The formal statement permits \(a=0\) or \(m=0\) as well. We have stated the
positive case needed for a nonempty family of compatible powers; no such
corner case is used to obtain a stronger arithmetic conclusion.
The threshold is uniform in the polynomials, the points, their degrees, and
the finite sets of coordinates and monic factors subject to the displayed bounds.


\section{Determinants, relation modules, and denominators}\label{sec:denominators}
Two finite-dimensional arguments separate the construction of Selmer classes
from the estimates eventually needed for volumes.

\subsection{A sharp determinant in a family of relations}
\begin{proposition}[determinant specialization]\label{prop:fitting}
Let \(\mathcal O\) be a discrete valuation ring, let \(c\in\mathcal O\setminus\{0\}\),
and let \(A(X)\) be an \(s\)-by-\(d\) matrix over
\(\mathcal O[X_1,\ldots,X_q]\).
For a homomorphism \(e:\mathcal O\to\mathcal O'\) into a discrete valuation ring
with \(e(c)\ne0\), and \(u\in(\mathcal O')^q\), let
\[
 \mathcal S_{e,u}
 =\ker\bigl(e(A)(u):(\mathcal O'/e(c))^d
                         \longrightarrow(\mathcal O'/e(c))^s\bigr).
\]
Normalize the discrete valuation \(v'\) by giving a uniformizer value \(1\).
Fix \(u_0\in\mathcal O^q\). There is a polynomial
\(F\in\mathcal O[X_1,\ldots,X_q]\) such that
\[
 v(F(u_0))=\length_{\mathcal O}\mathcal S_{\mathrm{id},u_0},
 \qquad
 \length_{\mathcal O'}\mathcal S_{e,u}
                  \le v'(e(F)(u))
\]
for every such \(e,u\), with \(v'(0)=+\infty\).
\end{proposition}
\begin{proof}
Append \(cI_d\) as additional rows of \(A(X)\), and call the resulting
\((s+d)\)-by-\(d\) matrix \(R(X)\). At any specialization into a discrete
valuation ring, elementary row and column operations put \(A\) in Smith
normal form, with nonzero diagonal entries of valuations
\(a_1,\ldots,a_t\), and \(d-t\) zero directions.
If \(h=v'(e(c))\), the kernel modulo \(e(c)\) has length
\[
 \sum_{i=1}^t\min(a_i,h)+(d-t)h .
\]
This follows coordinate by coordinate: the solutions of
\(\pi^{a_i}x=0\) modulo \(\pi^h\) form a module of length \(\min(a_i,h)\);
a zero direction has length \(h\).

The ideal generated by the maximal minors of \(R\) has exactly that valuation.
One way to see this is to compute its cokernel before taking minors.
The row lattice of \(R\), after Smith operations on \(A\), is generated
in coordinate \(i\) by \(\pi^{a_i}\) and \(\pi^h\), or by \(\pi^h\) alone
in a zero direction. It is a full lattice of index valuation equal to the
displayed sum. A basis matrix for a full lattice has determinant of valuation
its index, and elementary row operations show that its determinant ideal is
the ideal of maximal minors of any generating matrix.
In a discrete valuation ring a finitely generated ideal has valuation equal
to the minimum of the valuations of its generators. Thus the kernel length
is the minimum valuation among these minors.

At \(u_0\), choose a maximal minor attaining this minimum, and let \(F(X)\)
be the determinant of those same rows of \(R(X)\).
At every other specialization its valuation is at least the minimum there.
This proves both claims, including the case when that particular determinant
specializes to zero. The appended minor \(c^d\) ensures that the minimum is
always finite.
\end{proof}

In the Selmer argument, a separate finite-coordinate theorem presents the
cocycles and their local conditions by such relations.
The proposition then supplies one polynomial that bounds every specialized
length and is sharp at the selected base point. Neither this determinant
argument nor Theorem~\ref{thm:many-torsion} proves the finite-coordinate theorem
or supplies the base point with short Selmer group.
\src{New/Selmer/Corepresentable.lean}{finite relation coordinates}
\src{New/Selmer/LinePointPowers.lean}{exists_sharp_fitting_poly}

\subsection{Simultaneously non-small linear forms}
\begin{lemma}\label{lem:generic-residue}
Let \(K=\overline{\Q}_\ell\), and let \(L_1,\ldots,L_d\) be nonzero linear
forms on \(K^s\). There is \(z\in\mathcal O_K^s\) such that
\[
 \abs{L_k(z)}=\max_p\abs{\operatorname{coeff}_p L_k}\qquad(1\le k\le d).
\]
\end{lemma}
\begin{proof}
Divide each form by one of its coefficients of greatest absolute value.
All the resulting coefficients are integral, and each form has nonzero
reduction in the residue field. This residue field is infinite:
it contains a copy of every finite extension of \(\F_\ell\), obtained by
lifting an irreducible separable polynomial and taking a root in
\(\overline{\Q}_\ell\) in the corresponding unramified extension.
A finite union of proper linear hyperplanes cannot cover an affine space
over an infinite field. Explicitly, on the curve
\((1,t,\ldots,t^{s-1})\) each nonzero reduced linear form is a nonzero
polynomial in \(t\), with finitely many roots. Choose \(t\) avoiding all
of them, and lift its coordinates to \(\mathcal O_K\).
Every normalized form then evaluates to a unit. Multiplying back its
normalizing coefficient gives the assertion. The case \(s=0\) has no
nonzero linear forms and no substantive assertion.
\end{proof}

\begin{proposition}[scalar operators control a determinant]\label{prop:operators}
Fix \(\iota:\C\simeq\overline{\Q}_\ell\), and give complex numbers the
transported absolute value \(\abs z_\iota=\abs{\iota(z)}_\ell\).
Let \(Z\) be a \(d\)-dimensional complex subspace of \(V\), spanned by a set
\(\mathcal A\). Let \(j:V\to W\) be linear, \(x_1,\ldots,x_d\in Z\),
and \(\eta_1,\ldots,\eta_d\in V^*\). Let \(M\in\overline{\Q}_\ell\) and
\(\gamma\in\mathbb R\).
Suppose that for each nonzero \(p\in\mathcal A\) there is a linear map
\(G_p:V\to W\) with the following properties:
\begin{enumerate}[label=(\roman*)]
\item for every \(v\in\mathcal A\) some \(\lambda\in\C\) satisfies
      \(G_pv=\lambda jv\);
\item every \(\lambda\) with \(G_pp=\lambda jp\) satisfies
      \(\abs{M\,\iota(\lambda)}_\ell\ge e^{-\gamma}\);
\item if \(w\in Z\) and \(jw=G_p x_b\), then
      \(\abs{M\,\iota(\eta_a(w))}_\ell\le1\) for every \(a,b\).
\end{enumerate}
Then
\[
 \abs{\iota\bigl(\det(\eta_a(x_b))_{a,b}\bigr)}_\ell\le e^{d\gamma}.
\]
\end{proposition}
\begin{proof}
For \(d=0\) the determinant is \(1\), so suppose \(d>0\).
Choose a basis \(\omega_1,\ldots,\omega_d\) from \(\mathcal A\).
For \(G_p=G_{\omega_p}\), choose scalars
\(\lambda_{p,k}\) with \(G_p\omega_k=\lambda_{p,k}j\omega_k\).
By (ii),
\(\abs{M\iota(\lambda_{k,k})}_\ell\ge e^{-\gamma}>0\).
In particular the linear forms
\(L_k(z)=\sum_p M\iota(\lambda_{p,k})z_p\) are nonzero.
Lemma~\ref{lem:generic-residue} gives \(z_p\in\mathcal O_K\) with
\[
 \Lambda_k:=\sum_p M\iota(\lambda_{p,k})z_p,\qquad
 \abs{\Lambda_k}_\ell\ge e^{-\gamma}\quad(1\le k\le d).
\]

Write \(x_b=\sum_k P_{bk}\omega_k\), and put
\(w_{p,b}=\sum_k P_{bk}\lambda_{p,k}\omega_k\).
Then \(w_{p,b}\in Z\) and \(jw_{p,b}=G_p x_b\).
Thus the matrix
\[
 Q_{ab}=\sum_p z_p M\iota(\eta_a(w_{p,b}))
\]
has integral entries by (iii). With \(E_{ak}=\iota(\eta_a(\omega_k))\)
we have \(Q=E\diag(\Lambda_k)\iota(P)^{\mathsf t}\).
Taking determinants, and using the same expression without the diagonal
matrix for \(\iota(\eta_a(x_b))\), gives
\[
 \det Q
 =\iota\bigl(\det(\eta_a(x_b))\bigr)\prod_k\Lambda_k.
\]
The left side has absolute value at most \(1\), while the product has
absolute value at least \(e^{-d\gamma}\). Division proves the result.
\end{proof}
\src{New/Selmer/DeterminantBound.lean}{norm_gram_det_le_of_operators}

The scalar operators need not be the same for different basis vectors.
The proposition combines them with integral coefficients chosen in the
local field, rather than assuming one initially clears every denominator.
It is valid without assuming that \(j\) is injective on all of \(V\);
the universal quantifier in (ii) is part of its hypothesis.

\subsection{The denominator statement needed for theta lifts}\label{sec:den-def}
Let \(O\) now denote an auxiliary pair of unitary spaces, not a valuation ring,
let \(\mathcal F(O)\) be its space of forms, and let \(e_0\) be an element
of the group algebra at its chosen finite places. For a subspace \(Z\) write
\(P_Z\) for its ordinary complex orthogonal projection.
The source's tested condition for a form \(y'\), relative to \(e_0,Z^\perp\),
means that, for every permissible split place \(w\) off the level,
\[
 y'\perp
 \bigl(\im e_0\cap\ker H_w(X_w)\cap Z^\perp\bigr).
\]
Here \(H_w\) is the Hecke-polynomial operator with the test value \(X_w\)
specified by the auxiliary pair. All these places, not a single convenient
one, occur in the condition.

For \(r_O\) the rank on the row side, the assertion
\(\operatorname{Den}_\iota(O,e_0,Z;\beta)\) is the following:
for every \(d=\dim Z\), and \(d\)-tuples \(y,y'\) of forms, if each \(y_a\)
is integral at every embedding into \(\C_\ell\), and each \(y'_b\) is integral
at the embedding attached to \(\iota\) and tested in the preceding sense,
then
\begin{equation}\label{eq:den}
 \abs{\iota\bigl(\det(\ip{y_a}{P_Zy'_b})_{a,b}\bigr)}_\ell
 \le \exp(d\,r_O\,\beta).
\end{equation}
The definition does not assert that \(P_Z\) preserves an integral lattice.
Its permitted denominators are controlled only after pairing with these
integral and tested vectors.
\src{New/Selmer/FitTwists.lean}{DenominatorLe}

For a twist with the source's short Selmer groups, the required value is
\(\beta=\beta_\ell\log\ell\), with \(\beta_\ell\) defined in
Section~\ref{sec:budget}.
The formal argument constructs the \(G_p\) of
Proposition~\ref{prop:operators} by Ribet's method and patching.
It uses \(M=G^3\), where \(G\) is the order of the chosen finite group, and
\(\gamma=r_O\beta_\ell\log\ell\).
The proposition then proves \eqref{eq:den}; it does not construct the
operators. That source-specific construction remains the input
in Section~\ref{sec:fit}.


\section{The mirror estimate for tied cocycles}\label{sec:mirror}
This is the finite algebra behind the reciprocity step in the short-Selmer
construction. We prove it without leaving the key identity to a formal
linear-combination command.

Let \(A\) be a commutative ring, \(I\) an ideal with \(A/I\) finite, and
\(M=(A/I)^n\) with a specified coordinate \(i\).
Let \(\Gamma\) be a group, \(c:\Gamma\to\Gamma\) an involutive automorphism,
\(\rho:\Gamma\to\mathrm{GL}_n(A)\), and
\(\chi:\Gamma\to A^\times\).
Let \(H\) be a normal \(c\)-stable subgroup of \(\Gamma\), contained in
a subgroup \(U\) that acts trivially on \(M\).
Fix \(\gamma\in\Gamma\) such that \([k,\gamma]\in H\) for every \(k\in\Gamma\);
our convention is \([a,b]=aba^{-1}b^{-1}\).
Suppose \(s-1,s+1\) are units in \(A\), and
\[
 \rho(\gamma)=sI_n,\qquad \rho(c\gamma)^{-1}=sI_n,\qquad
 \chi(\gamma)=\chi(c\gamma)=1.
\]
The scalar \(s\) is also a unit, since \(sI_n\) is invertible and \(n\ge1\).

Let \(y:\Gamma\to M\) be a \(1\)-cocycle with \(y(\gamma)=0\), and choose
\(g_0\) with \(y(g_0)_i\) a unit of \(A/I\).
Assume the \(A\)-span of the matrices \(\rho(\Gamma)\) is all \(M_n(A)\).
Choose \(z\in U\) with \(cz\in U\).
For a set \(W\), suppose each \(w\in W\) gives functions
\(x_w:\Gamma\to M^\vee\), where
\(M^\vee=\operatorname{Hom}_{\rm add}(M,\Q/\Z)\).
Their required equivariance is
\begin{equation}\label{eq:mirror-equivariance}
 x_w(a k b k^{-1})(v)
   =x_w(a)(v)+x_w(b)(\chi(k)\rho(k)^{-1}v)
 \quad(a,b\in U,\ k\in\Gamma,\ v\in M).
\end{equation}
Assume distinct \(w\)'s have distinct restrictions \(x_w|_U\), and
\begin{equation}\label{eq:mirror-isotropic}
 x_w(zh)(t\,y(zh))
       =x_w(c(zh))(t\,y(c(zh)))
 \quad(h\in H,\ t\in A).
\end{equation}
No \(A\)-linearity of \(x_w(g)\) is assumed: its additivity and the
quantifier over \(t\) are both needed.

\begin{theorem}[mirror estimate]\label{thm:mirror}
If \(g\in\Gamma\) and
\(\rho(g)=\alpha I_n,\ \rho(cg)^{-1}=\alpha'I_n\),
put \(\delta=\alpha-\chi(cg)\alpha'\in A\).
Then \(W\) is finite and
\[
 \#W\le\#M[\delta],\qquad
 M[\delta]=\{v\in M:\delta v=0\}.
\]
\end{theorem}
\begin{proof}
We suppress \(w\) temporarily.
Taking \(k=1\) in \eqref{eq:mirror-equivariance} shows that \(x|_U\)
is an additive homomorphism, in multiplicative group notation:
\(x(ab)=x(a)+x(b)\), and \(x(1)=0\).
Taking \(a=1\) gives
\[
 x(k b k^{-1})(v)=x(b)(\chi(k)\rho(k)^{-1}v)\quad(b\in U).
\]
On \(U\), the cocycle \(y\) is likewise additive, because \(U\) acts
trivially on \(M\).
For \(b\in H\), the general cocycle identity gives
\[
 y(k b k^{-1})=(1-\rho(k b k^{-1}))y(k)+\rho(k)y(b)
             =\rho(k)y(b),
\]
as \(H\) is normal and acts trivially on \(M\).

Put \(\tau=[g_0,\gamma]\in H\).
The cocycle identity, scalar action of \(\gamma\), and \(y(\gamma)=0\)
give \(y(\tau)=(1-s)y(g_0)\). Thus \(y(\tau)_i\) is a unit.
The vectors \(t\rho(k)y(\tau)\), \(t\in A,k\in\Gamma\), span \(M\)
additively. Indeed the matrices \(\rho(k)\) span \(M_n(A)\).
Every vector in \(M\) is a matrix over \(A/I\) applied to \(y(\tau)\):
put the desired coordinates, divided by its unit \(i\)-th coordinate,
in column \(i\), and set the other columns to zero.
Lift that matrix to \(M_n(A)\) and use the spanning assumption.

\emph{Polarization.}
For \(a\in U\) with \(ca\in U\), define
\(q_t(a)=x(a)(t y(a))-x(ca)(t y(ca))\).
By \eqref{eq:mirror-isotropic}, \(q_t(zh)=0\) for \(h\in H\).
Expand
\[
 q_t(zab)-q_t(za)-q_t(zb)+q_t(z)=0\quad(a,b\in H).
\]
Additivity of \(x,y\) on \(U\) and of their conjugates gives
\begin{equation}\label{eq:mirror-polar}
 x(a)(t y(b))+x(b)(t y(a))
  =x(ca)(t y(cb))+x(cb)(t y(ca)).
\end{equation}

Fix \(b\in H\) and \(k\in\Gamma\), and insert
\(a=k\tau k^{-1}\) into \eqref{eq:mirror-polar}.
As additive functions of \(t\in A\), put
\begin{align*}
 A_0(t)&=x(b)(t\rho(k)y(\tau)),\\
 B_0(t)&=x(\tau)(\chi(k)t\rho(k)^{-1}y(b)),\\
 C_0(t)&=x(cb)(t\rho(ck)y(c\tau)),\\
 D_0(t)&=x(c\tau)(\chi(ck)t\rho(ck)^{-1}y(cb)).
\end{align*}
Equation~\eqref{eq:mirror-polar} says \(A_0+B_0=C_0+D_0\).
Replace \(b\) by \(\gamma b\gamma^{-1}\), which still lies in \(H\).
Equivariance and the scalar hypotheses give
\begin{align*}
 x(\gamma b\gamma^{-1})(v)&=x(b)(s^{-1}v),
 &y(\gamma b\gamma^{-1})&=s y(b),\\
 x(c(\gamma b\gamma^{-1}))(v)&=x(cb)(s v),
 &y(c(\gamma b\gamma^{-1}))&=s^{-1}y(cb).
\end{align*}
Thus the new polarization equation is
\[
 A_0(s^{-1}t)+B_0(st)=C_0(st)+D_0(s^{-1}t).
\]
Subtract the original equation evaluated at \(s^{-1}t\).
It gives
\((B_0-C_0)((s-s^{-1})t)=0\).
Since \(s-s^{-1}=s^{-1}(s-1)(s+1)\) is a unit, \(B_0=C_0\);
hence \(A_0=D_0\).
We have proved the key identity
\begin{equation}\label{eq:mirror-key}
 x(b)(t\rho(k)y(\tau))
   =x(c\tau)(\chi(ck)t\rho(ck)^{-1}y(cb))
 \quad(b\in H,\ t\in A,\ k\in\Gamma).
\end{equation}

\emph{The scalar difference.}
Let \(\ell_w=x_w(c\tau)\).
Apply \eqref{eq:mirror-key} with \(b=c\tau\):
\[
 \ell_w(t\rho(k)y(\tau))
   =\ell_w(\chi(ck)t\rho(ck)^{-1}y(\tau)).
\]
Compare this identity at \(kg,t\) with its instance at
\(k,\chi(cg)\alpha't\). The right sides agree, while the difference of
the left sides is
\(\ell_w(\delta t\rho(k)y(\tau))\).
The spanning assertion above shows that \(\ell_w\) annihilates
\(\delta M\). It therefore descends to a character of \(M/\delta M\).

\emph{Injectivity.}
Suppose \(\ell_w=\ell_v\).
Subtract \eqref{eq:mirror-key} for \(w,v\). Its right side vanishes.
Spanning gives \(x_w(b)=x_v(b)\) for every \(b\in H\).
For \(a\in U\), equivariance gives
\[
 x_w([a,\gamma])(m)
   =x_w(a)(m)-x_w(a)(s^{-1}m)
   =x_w(a)((1-s^{-1})m).
\]
The same formula holds for \(v\), and the commutator belongs to \(H\).
As \(1-s^{-1}\) is a unit, multiplication by it is onto \(M\).
Thus \(x_w(a)=x_v(a)\) for all \(a\in U\), whence \(w=v\).

We have an injection
\(W\hookrightarrow\operatorname{Hom}_{\rm add}(M/\delta M,\Q/\Z)\).
A finite abelian group has as many such characters as elements:
decompose it into cyclic groups, and note that a cyclic group of order
\(r\) has exactly \(r\) possible images of a generator in \(\Q/\Z\).
Finally a homomorphism \(\delta:M\to M\) of a finite group has
\(\#(M/\delta M)=\#\ker\delta\), by counting fibers and image.
This proves finiteness and the asserted bound.
\end{proof}
\src{New/Selmer/LinePointPowers.lean}{card_le_torsion_of_isotropic}

The formal statement writes \(\rho\) as a homomorphism to the multiplicative
monoid of matrices, and \(\chi\) to the multiplicative monoid of \(A\).
Since their domain is a group, their images are units; the formulation
above is the same condition. Stability \(c(H)\subseteq H\), together with
\(c^2=1\), is equality. The proof does not add a condition \(c(U)=U\).

The next two sections prove the global isotropy and the bounded-pencil
construction. Balancing the particular automorphic rows is a separate
local arithmetic calculation.



\section{Global Selmer groups and bounded pencils}\label{sec:global-selmer}
We now connect the finite mirror estimate to global reciprocity.
The conventions at conjugate places matter. Write \(G=G_E\), let
\(c(g)\) denote conjugation by a fixed complex conjugation of the
absolute Galois group of the maximal real subfield, and put
\(g^-=(c(g))^{-1}\), \(g^+=g\).
Thus \(c^2=1\), and the cyclotomic character satisfies
\(\chi(c(g))=\chi(g)\). A tame generator at the conjugate place is
\(\tau^-=(c(\tau))^{-1}\), rather than \(c(\tau)\).

Throughout this section \(E\) is a CM number field,
\(f=[E:\Q]/2\), and \(A\) is the ring of integers of a finite
extension of \(\Q_\ell\). Its valuation \(v_A\) has value \(1\) on
a uniformizer \(\pi\), and its residue field has order \(q_A\).
Lengths in this section are unnormalized \(A\)-lengths.
For nonzero \(a\in A\) put \(M_a=(A/a)^n\).
If \(C\) is a finite \(A\)-module, then
\(\#C=q_A^{\length_A C}\); this follows by a composition series,
whose factors are all the residue field.

A relation is a finitely supported family \(R=(R_g)_{g\in G}\)
of matrices over \(A\); it requires
\(\sum_gR_gz(g)=0\).
For a matrix-valued function \(r:G\to M_n(A)\), a finite set \(S\)
of finite places, and any set \(\mathcal L\) of relations, let
\(\mathcal S(r,S,\mathcal L;B)\) consist of the locally constant
functions \(z:G\to B^n\) satisfying
\[
 z(gh)=z(g)+r(g)z(h),\qquad
 z|_{I_v}=0\ (v\notin S),\qquad
 \sum_gR_gz(g)=0\ (R\in\mathcal L).
\]
Vanishing outside \(S\) is required at every place above \(v\).
This definition makes sense even when \(r\) is not a representation.
That flexibility is essential for polynomial specialization.

\subsection{Finite coordinates, before specializing}
\begin{proposition}[a uniform finite presentation]\label{prop:selmer-coordinates}
Let \(\Lambda\) be a commutative ring, \(m\ge1\), and
\(r:G\to M_n(\Lambda)\) any function.
Fix a neighborhood \(W\) of \(1\), a finite set \(S\), and relations
\(\mathcal L\).
There are finitely many coordinates \(g_j,i_j\) and a submodule
\(R\subseteq\Lambda^J\), containing \(m\Lambda^J\), such that
evaluation identifies \(\mathcal S(r,S,\mathcal L;B)\) with
\[
 \{b\in B^J:\textstyle\sum_j u_jb_j=0\text{ for every }u\in R\}
\]
for every \(\Lambda\)-module \(B\) killed by \(m\) on which
\(r(g)\) acts as the identity for all \(g\in W\).
If \(\Lambda\) is Noetherian, \(R\) is finitely generated.
\end{proposition}
\begin{proof}
Choose an open normal subgroup \(U\subset W\), and let \(F\) be
its fixed field. The restriction to \(U\) of any allowed \(z\) is
an additive homomorphism, because \(r(u)\) acts trivially on \(B^n\).
It is unramified outside the finitely many places of \(F\)
over \(S\).

There are only finitely many continuous homomorphisms
\(U\to\Z/m\) unramified outside these places.
Here is the class field theory justification, including the
bounded-conductor point. At a nonarchimedean completion of \(F\)
the subgroup of \(m\)-th powers in its multiplicative group is
open. For residue characteristic prime to \(m\), this follows
from Hensel's lemma on the principal units. In residue
characteristic dividing \(m\), a sufficiently deep principal-unit
group is carried by the convergent local logarithm to an ideal
in the local field; multiplication by \(m\) has open image in
that ideal. Consequently every character of exponent \(m\)
has bounded conductor at each of the permitted places.
Global reciprocity now identifies the characters with characters
of one finite ray class group. At the real places there is
only the finite sign choice.

Intersect their kernels and take the core in \(G\).
The result is an open normal subgroup \(V\subset U\).
Every allowed \(z\) vanishes on \(V\).
Indeed additive characters into \(\Q/\Z\) separate points of
every abelian group; for an \(m\)-torsion group their values lie
in \((\Q/\Z)[m]\simeq\Z/m\).
Applying such a character to each coordinate of \(z|_U\)
gives one of the homomorphisms just killed.
For clarity, the separation assertion follows by mapping a
cyclic subgroup generated by a nonzero element of order \(d\)
to \(1/d\pmod\Z\), and extending the character to the ambient
abelian group. Such extension holds because \(\Q/\Z\) is
divisible: when adjoining one new generator, its prescribed
multiple has a preimage; a maximal extension then has full
domain.

Since \(z(gv)=z(g)\) for \(v\in V\), choose representatives for
the finite set \(G/V\). Use their \(n\) coordinates as \(J\).
Conversely a tuple \(b\in B^J\) defines a function constant on
these cosets. In the free module \(\Lambda^J\), impose all its
cocycle equations, its inertia equations outside \(S\), the
relations \(\mathcal L\), and \(m\Lambda^J\).
Let \(R\) be the submodule generated by these equations.
Their simultaneous annihilator is exactly the displayed set,
and a function on a finite quotient is locally constant.
This proves both directions, with no representation assumption
on \(r\). Finite generation of \(R\) is the Noetherian property.
\end{proof}
\src{New/Selmer/Corepresentable.lean}{exists_corepresenting_coords}

Apply the proposition with
\(\Lambda=A[X_v:v\in J]\) and
\[
 r_y(g)=\prod_{v\in J}y_v^{\widetilde h_v(g)}\rho(g),
 \qquad 0\le\widetilde h_v(g)<\ell^N.
\]
A neighborhood on which every \(h_v\) vanishes and
\(\rho\bmod\ell^M\) is the identity works uniformly.
It follows that the finite relation matrices, and hence
Proposition~\ref{prop:fitting}, apply before the \(y_v\) are
roots of unity. In particular, one may choose a determinant
polynomial sharp at any specified tuple; no continuity
assertion about a nonexistent characteristic-zero
representation at a general tuple is needed.

\subsection{Auxiliary anticyclotomic primes}
We use the following precise consequence of global class field
theory and Chebotarev's theorem. A character is written
additively in the statement.
\begin{proposition}\label{prop:auxiliary-prime}
Let \(T_0,R\) be finite sets of finite places of \(E\).
Let \(z\) belong to the normal subgroup generated by inertia at
\[
 R\setminus(T_0\cup\bar T_0\cup\{v\mid\ell\}).
\]
For every neighborhood \(W\) of \(z\) and every \(N\ge1\),
there are a place \(P\notin R\cup\bar R\) of rational prime
norm, a place above it, an arithmetic Frobenius \(\phi\in W\),
a tame generator \(\tau\), and a character
\(h:G\to\Z/\ell^N\) such that
\[
 \begin{gathered}
 NP\equiv1\pmod{\ell^N},\qquad h(\tau)=1,\quad h(\phi)=0,\\
 h(c(g))=-h(g),\qquad
 h|_{D_v}=0\quad(v\in T_0\cup\bar T_0\cup\{v\mid\ell\}).
 \end{gathered}
\]
The character is tame and unramified outside \(P,\bar P\).
It is induced by an idele-class character evaluated on
\(x/\bar x\).
One may enlarge the finite excluded set while retaining the
assertion, provided the displayed inertia condition on \(z\)
still holds.
\end{proposition}
\begin{proof}
We give the construction, using the usual reciprocity
isomorphism between finite abelian extensions and finite
idele-class characters. Choose a finite set \(T\), containing
\(T_0,\bar T_0\) and the places over \(\ell\), whose ideal classes
generate the class group. The added places can avoid the
places at which inertia is allowed in the displayed condition.
To see this, each ideal class contains a prime outside any
fixed finite set, by Chebotarev applied to the Hilbert class
field. Thus one may choose one representative prime in each
of a finite set of generating classes.

The group of \(T\)-units is finitely generated. Let \(u_1,\ldots,u_s\)
be generators, including its torsion generators, and form
\[
 L=E(\mu_{\ell^N},u_1^{1/\ell^N},\ldots,u_s^{1/\ell^N}).
\]
Take a finite Galois closure if needed.
At a place outside \(T\) of residue characteristic different
from \(\ell\), adjoining an \(\ell^N\)-th root of a unit is
unramified: after adjoining the necessary residue roots its
defining polynomial has unit derivative, so Hensel's lemma
lifts it in an unramified extension.
The same holds for the cyclotomic extension.
Consequently the normal subgroup in the hypothesis fixes
\(L\), and \(z\) fixes \(L\).

Choose an open normal subgroup \(V\) so that \(zV\subset W\),
and apply Chebotarev to the finite Galois extension cut out by
\(V\) together with \(L\). We use its degree-one-prime form:
an element of \(G_E\) can be approximated, with a chosen
place above the prime, by Frobenius elements at primes of
\(E\) of rational prime norm. This form follows by applying
Chebotarev to a Galois closure over \(\Q\), extending the
specified automorphism which fixes \(E\); its fixed embedding
of \(E\) gives a degree-one prime. It gives infinitely many
choices, so all finitely many excluded places and the
ramified places of \(V\) can be avoided.
We obtain an arithmetic Frobenius \(\sigma\in zV\) at \(P\)
fixing \(L\). In particular \(NP\equiv1\pmod{\ell^N}\), and
every \(u_j\) reduces to an \(\ell^N\)-th power in the residue
field at \(P\).

Choose a surjection
\(\eta:(\mathcal O_E/P)^\times\to\Z/\ell^N\), normalized to
take the chosen tame generator to \(1\) under local
reciprocity. It kills the reductions of all \(T\)-units.
The group of ideles which are units outside \(T\) surjects
onto the idele class group: the class-group generating
property of \(T\) permits the valuations of an arbitrary
idele outside \(T\) to be canceled by a principal idele.
On this group, reduction of its \(P\)-component followed by
\(\eta\) kills the principal ideles in the group, precisely
because they are the \(T\)-units. It therefore defines a
continuous character \(\eta_0\) of the idele class group.
It is trivial on the full local multiplicative groups at
\(T\), on units at every place outside \(T\cup\{P\}\), and
on principal units at \(P\).

Set \(\eta_1(x)=\eta_0(x/\bar x)\), and let \(h\) be the
corresponding Galois character. Because \(P\) has degree one,
\(NP\equiv1\pmod{\ell^N}\), and the finitely many ramified
places were avoided, \(P\ne\bar P\).
The construction gives all the asserted local conditions,
and \(h(\tau)=1\).
Finally replace \(\sigma\) by
\(\phi=\sigma\tau^{-h(\sigma)}\).
It is still a Frobenius lift. The chosen \(P\) is unramified
in the quotient \(G/V\), so this replacement remains in
\(zV\subset W\). It has \(h(\phi)=0\).
\end{proof}
\src{New/Selmer/Corepresentable.lean}{exists_isAnticyclotomicAt}

\subsection{Balanced rows and pencils}
Here is the full numerical condition needed for the reciprocity
argument. For each place \(v\) choose a decomposition group
\(D_v\subset G\). Put \(C_v(\rho,\mathcal L;a)\) equal to the
image on \(D_v\) of all global locally constant cocycles
with values in \(M_a\) satisfying \(\mathcal L\), without
a ramification restriction. It is an image of cocycles,
not an arbitrarily chosen local crystalline subspace.
Its length can depend on the chosen place above \(v\);
write \(\min_{D_v}\) for the least such length.

\begin{definition}\label{def:balanced}
The row \((\rho,\mathcal L,\gamma)\) is balanced up to
\(\beta\) modulo \(\ell^M\) if the following hold.
The reduction of \(\rho\) modulo \(\ell^M\) is continuous,
\(\rho\) is unramified outside \(S\), and its matrices span
\(M_n(A)\). There is \(s\in A^\times\) with \(s-1,s+1\)
units such that \(\gamma\) is inertial at a place of
\(S\) away from \(\ell\),
\[
 \rho(\gamma)=\rho(\gamma^-)=sI.
\]
The residue field has more than \(4n+2\) elements.
Every coboundary \(g\mapsto(\rho(g)-1)b\) satisfies
\(\mathcal L\), for every coefficient quotient \(A/a\).
Finally there are \(g,\alpha,\alpha'\) with
\(\rho(g)=\alpha I,\rho(g^-)=\alpha'I\) such that,
for every nonzero \(a\mid\ell^M\),
\begin{equation}\label{eq:balanced-length}
 \sum_{v\in S}\min_{D_v}\length_A C_v(\rho,\mathcal L;a)
 +n v_A(\alpha-\chi(c(g))\alpha')
 \le (|S|+f)\length_A M_a+\beta .
\end{equation}
\end{definition}
The minimum exists: it is the minimum of a nonempty set
of nonnegative integers. A bound \(\beta=+\infty\) is
vacuous; all uses below have finite \(\beta\).

Enlarge the fixed forbidden set \(S_0\supset S\) to contain
its conjugate and the places over \(\ell\).
A pencil consists of a finite auxiliary set \(J\), a helper
\(P_0\in J\), data \((h_v,\tau_v,\phi_v)\) from
Proposition~\ref{prop:auxiliary-prime}, and slopes
\(t_{v,i}\in A\) for finitely many rows \(i\).
The sets \(J,\bar J,S_0\) are pairwise disjoint where
required, and every \(v\in J\) has rational prime norm.
Its defining compatibilities are
\begin{equation}\label{eq:pencil}
 \begin{gathered}
 h_w(\tau_v)=\delta_{wv},\qquad h_v(\phi_v)=0,\\
 h_w(\phi_v)=\delta_{wP_0}\quad(v\ne P_0),\qquad
 \phi_{P_0}\in W,\\
 t_{v,i}\ne0\ \Longrightarrow\
       \rho_i(\phi_v^\pm)\equiv1\pmod{\ell^M}.
 \end{gathered}
\end{equation}
The characters are trivial on decomposition groups over
\(\ell\). The first line for \(w\ne v\) follows from
unramifiedness; the second line is an extra prescription
at old Frobenii, and must not be omitted.
The ties on row \(i\) are
\(z(\tau_v^\pm)=t_{v,i}z(\phi_v^\pm)\).
Take the helper slope to be zero initially.

For one row suppress \(i\), and define
\[
 X_a(T)=\mathcal S\bigl(\rho,S\cup T\cup\bar T,
       \mathcal L\cup\{\text{ties at }J\}\cup\{z(\gamma)=0\};
       A/a\bigr).
\]
Thus \(X_a(J\cup\{P\})\) has the old ties only;
the new prime is relaxed.

\subsection{Global reciprocity detects an unseen class}
The classical background used here consists of local Tate
duality, global reciprocity for the cup product, and the
Greenberg--Wiles cardinal formula. We state the latter
explicitly to specify every factor. For a finite Galois
module \(M_0\), its Cartier dual
\(M_0^*=\operatorname{Hom}(M_0,\Q/\Z(1))\), and local
cohomology subgroups \(L_v\) unramified at almost every
finite place, the formula is
\begin{equation}\label{eq:GW}
 \frac{|\operatorname{Sel}_{L}(M_0)|}
      {|\operatorname{Sel}_{L^\perp}(M_0^*)|}
 =\frac{|H^0(E,M_0)|}{|H^0(E,M_0^*)|}
      \prod_v\frac{|L_v|}{|H^0(E_v,M_0)|}.
\end{equation}
Here \(L_v^\perp\) is the annihilator for the perfect local
cup-product pairing. At a complex place the factor is
\(|M_0|^{-1}\). These standard duality theorems are used
as established background, not as new results of this
exposition; see \cite[Theorem 1.8]{smithGW} and \cite{milneADT}.
The compact source obtains the inequality needed below
by its equivalent split-prime relation construction.
Using \eqref{eq:GW} avoids carrying those additional
relation coordinates through the prose.

\begin{theorem}[a group which no new prime detects is short]
\label{thm:unseen}
Suppose the row is balanced as in
Definition~\ref{def:balanced}, \(a\ne0\) divides \(\ell^M\),
and \(X_a(J)\) contains a cocycle \(y\) with a unit
coordinate at some group element.
Suppose also that for every new auxiliary datum
\((P,h_P,\tau_P,\phi_P)\) which extends the pencil for
every slope on this row,
\begin{equation}\label{eq:unseen-inclusion}
 \{(z(\phi_P),z(\phi_P^-)):z\in X_a(J)\}
 \subseteq
 \{(z'(\tau_P),z'(\tau_P^-)):
                      z'\in X_a(J\cup\{P\})\}.
\end{equation}
Then \(\length_A X_a(J)\le\beta\).
\end{theorem}
\begin{proof}
Put \(M_0=M_a\) and \(S'=S\cup J\cup\bar J\).
At each \(v\in S'\), use the local condition \(C_v\) obtained
by restricting the global cocycles obeying \(\mathcal L\)
and the old ties, and then quotient by local coboundaries
to obtain \(L_v\subset H^1(E_v,M_0)\).
These conditions contain the local coboundaries.
Indeed global coboundaries satisfy \(\mathcal L\).
At a new tame generator \(\rho(\tau_v^\pm)=1\);
if its slope is nonzero,
\(\rho(\phi_v^\pm)\equiv1\pmod{\ell^M}\);
if its slope is zero, no Frobenius term occurs.
This checks the ties for every coboundary.
Choose the place above each \(v\in S\) minimizing the
length without the ties; at \(J,\bar J\) use the places
specified by the pencil. Adding ties can only shrink
a local image.

Let \(Y=\operatorname{Sel}_L(M_0)\) and
\(W^*=\operatorname{Sel}_{L^\perp}(M_0^*)\).
Both are finite. Multiplication by \(s-1\) is invertible,
so \(H^0(E,M_0)=0\). On the Cartier dual, \(\gamma\)
acts by \(s^{-1}\), since \(\chi(\gamma)=1\);
hence its global invariants also vanish.
Each local coboundary group has cardinality
\(|M_0|/|H^0(E_v,M_0)|\). Substituting
\(|C_v|=|L_v|\,|M_0|/|H^0(E_v,M_0)|\)
in \eqref{eq:GW} gives the exact identity
\begin{equation}\label{eq:GW-cocycles}
 |Y|\,|M_0|^{|S'|+f}=|W^*|\prod_{v\in S'}|C_v|.
\end{equation}
Outside \(S'\), the unramified local condition contributes
factor \(1\). If necessary include the places over \(\ell\)
in the set used to apply duality; this does not change the
identity because those outside \(S'\) still have their
unramified condition.

Every cohomology class has a unique cocycle representative
vanishing at \(\gamma\): evaluation of a coboundary at
\(\gamma\) is \((s-1)b\), a bijection on \(M_0\).
Consequently \(X_a(J)\) injects into \(Y\).
It need not equal \(Y\); global relations may be stronger
than their separate local images.

We next bound \(W^*\) by the mirror theorem. Let
\[
 U=\ker(\rho\bmod\ell^M)\cap
                 \ker(\chi\bmod\ell^{M+1}).
\]
Restriction of a class of \(W^*\) to \(U\) is a homomorphism
\(x_w:U\to M_0^\vee\), after fixing the compatible roots of
unity to identify the target with \(\Q/\Z\).
Distinct classes give distinct homomorphisms.
Here is the requisite injectivity argument.
If a class restricts to zero, inflation--restriction
places it in \(H^1(G/U,M_0^*)\).
The image of \(\gamma\) is central in \(G/U\), because
its representation is scalar and its cyclotomic value is
\(1\). For a cocycle \(b\) of this quotient, commuting
\(\gamma\) and \(g\) gives
\[
 (\gamma-1)b(g)=(g-1)b(\gamma).
\]
The operator \(\gamma-1=s^{-1}-1\) is invertible on the
dual, so \(b(g)=(g-1)(\gamma-1)^{-1}b(\gamma)\) is a
coboundary. This proves injectivity.
The cocycle formula on the dual also gives exactly
\eqref{eq:mirror-equivariance}; this formula concerns
restriction to \(U\), so it is independent of the
representatives chosen for \(w\).

Let \(H\) be the intersection with \(U\) and \(c(U)\) of
the commutator of \(G\) with the normal closure of
\(\{\gamma,\gamma^-\}\). This is normal and \(c\)-stable.
All \([k,\gamma]\) and their conjugates lie in \(U\cap c(U)\):
their matrices are the identity since the matrices of
\(\gamma,\gamma^-\) are scalar, and their cyclotomic
characters are \(1\). Thus \(H\) has the commutators
required in Theorem~\ref{thm:mirror}.
Put \(z=\tau_{P_0}\). Both \(z\) and \(c(z)\) lie in \(U\)
because the row is unramified at the helper prime and
inertia away from \(\ell\) fixes the cyclotomic tower.

Fix \(h\in H\). The element \(zh\) belongs to the normal
closure of inertia at \(S_0\cup\{P_0,\bar P_0\}\), away
from \(\ell\) and the old nonhelper primes.
For every old character,
\[
 h_w(zh)=\delta_{wP_0};
\]
characters kill commutators. Apply
Proposition~\ref{prop:auxiliary-prime} with the old
nonhelper primes among \(T_0\).
Choose its Frobenius \(\phi_P\) sufficiently close to
\(zh\) that \(\phi_P,c(\phi_P)\in U\), that all old
characters take these prescribed values, and that all
the finitely many functions \(y,x_w\) take the same
values as at \(zh,c(zh)\).
Avoid also every ramified place of the finitely many
dual cocycles. Its new character kills the old
nonhelper decomposition groups and itself vanishes
at \(\phi_P\). Thus \eqref{eq:pencil} holds with any
new slope on the chosen row, and with zero slopes
on the other rows.

By \eqref{eq:unseen-inclusion} there is
\(y'\in X_a(J\cup\{P\})\) with
\(y'(\tau_P^\pm)=y(\phi_P^\pm)\).
For every \(t\in A\), apply global cup-product reciprocity
to \(ty'\) and a representative of \(w\in W^*\).
At old places the local pairing is zero by the
orthogonal local conditions. At all other finite
places except \(P,\bar P\), both classes are unramified,
so their pairing is zero. At a complex place it is zero.
At \(P\) and \(\bar P\) the dual cocycle is unramified
and \(\rho,\chi\) act trivially on the relevant
finite quotient. The unramified--tame local pairing is
evaluation of its Frobenius value on the tame value
of the primal cocycle. The tame generator at \(\bar P\)
is \((c(\tau_P))^{-1}\), which accounts for the minus
sign. Reciprocity therefore reads
\[
 x_w(\phi_P)(t y'(\tau_P))
       -x_w(c(\phi_P))(t y'(c(\tau_P)))=0.
\]
The identities with inverses introduce no extra
matrix factor here, since the indicated elements
act trivially on \(M_0\).
Replace the tame values by those of \(y\), and then
use the prescribed closeness to \(zh\).
We obtain \eqref{eq:mirror-isotropic}, for all \(w,h,t\).
Only finitely many \(w\) and \(t\bmod a\) occur, so
one neighborhood suffices.

Theorem~\ref{thm:mirror}, with
\(\delta=\alpha-\chi(c(g))\alpha'\), now gives
\[
 |W^*|\le|M_0[\delta]|,\qquad
 \length_A M_0[\delta]\le n v_A(\delta).
\]
For the latter bound use the cyclic summands of \(A/a\):
the kernel length in one coordinate is
\(\min(v_A(a),v_A(\delta))\).

At each place in \(J\cup\bar J\), an allowed local
cocycle is determined by its Frobenius value.
Its tame value is fixed by the tie; wild inertia
has no nontrivial homomorphism to an \(\ell\)-group
at this residue characteristic, and the maximal
pro-\(\ell\) tame inertia is generated by \(\tau_v^\pm\).
Hence \(|C_v|\le|M_0|\).
At \(S\) its length is at most the chosen minimum
without ties. Taking \(q_A\)-logarithms in
\eqref{eq:GW-cocycles}, canceling the \(2|J|\)
new-prime factors, and using the injection into \(Y\),
gives
\[
 \length_A X_a(J)
 \le n v_A(\delta)+
       \sum_{v\in S}\min_{D_v}\length_A
                    C_v(\rho,\mathcal L;a)
          -(|S|+f)\length_A M_a
 \le\beta .
\]
This includes the local bookkeeping explicitly.
\end{proof}
\src{New/Selmer/TiltingLemma.lean}{IsPencil.length_le_of_unseen_class}

\subsection{Tilting a local condition}
\begin{lemma}[generic slopes have short kernels]\label{lem:tilting}
Let \(I\ne0\) be an ideal of a DVR \(A\), let \(B\) be
an \(A\)-module, \(Y\subset B\), and
\(\Phi_0,\Phi_1:B\to(A/I)^d\) linear.
Outside at most \(d\) residues \(s\bmod\pi\),
\[
 \length_A(Y\cap\ker(s\Phi_0-\Phi_1))
 \le
 \length_A\bigl(Y\cap\ker\Phi_0
                   \cap\Phi_1^{-1}(\Phi_0Y)\bigr).
\]
\end{lemma}
\begin{proof}
Put \(Z=\sum_{s\in A}Y\cap\ker(s\Phi_0-\Phi_1)\).
For every summand \(\Phi_1z=s\Phi_0z\), so
\(\Phi_1Z\subset\Phi_0Z\).
Let \(P\subset A^d\) be the inverse image of \(\Phi_0Z\)
under reduction modulo \(I\). It is a free module of
rank at most \(d\).
Its map onto \(\Phi_0Z\) lifts through
\(\Phi_0:Z\to\Phi_0Z\), because \(P\) is free.
Choose such a lift \(\sigma:P\to Z\).
The map \(\Phi_1\sigma:P\to\Phi_0Z\) in turn lifts
through \(P\to\Phi_0Z\) to an endomorphism \(\Psi\)
of \(P\).
If \(\bar s\) is not a root of the characteristic
polynomial of \(\Psi\bmod\pi\), the map \(s-\Psi\)
is invertible on \(P\). There are at most
\(\operatorname{rank}P\le d\) exceptional residues.
For every other \(s\), reduction of
\((s-\Psi)P=P\) gives
\[
 (s\Phi_0-\Phi_1)Z\supset\Phi_0Z.
\]
If \(Z\) has finite length, the two kernel--image
exact sequences therefore imply
\[
 \length\ker((s\Phi_0-\Phi_1)|_Z)
       \le\length\ker(\Phi_0|_Z).
\]
If \(Z\) has infinite length, both kernels have
infinite length because their images are finite;
the same inequality in extended lengths still holds.
Every element of \(Z\cap\ker\Phi_0\) belongs to \(Y\)
and has \(\Phi_1z\in\Phi_0Z\subset\Phi_0Y\).
The asserted domain kernel is contained in the
first kernel above. These two containments prove
the claim.
\end{proof}
\src{New/Selmer/TiltingLemma.lean}{tilting}

\begin{theorem}[boundedly many tied primes]\label{thm:bounded-pencil}
For finitely many balanced rows over \(A\), a nonzero
\(a\mid\ell^M\), and a neighborhood \(W\) of an element
in the inertia subgroup specified in
Proposition~\ref{prop:auxiliary-prime}, there is an integer
\(J_0\), independent of \(N\), such that for every \(N\)
there is a pencil with \(|J|\le J_0\) and
\[
 \length_A X_{i,a}(J)\le\beta_i\quad\hbox{for all }i.
\]
\end{theorem}
\begin{proof}
Induct on \(v_A(a)\). For a unit \(a\) all modules
vanish. Proposition~\ref{prop:auxiliary-prime} supplies
the helper, and \(J=\{P_0\}\), with all slopes zero,
is a pencil.

Suppose the result holds at \(a\), and \(\pi a\mid\ell^M\).
Let \(J_1\) be its uniform bound.
For fixed \(N\), start with an inductive pencil.
Let \(T_i\) be the unramified-outside-\(S\) module
at modulus \(\pi a\), with relations \(\mathcal L_i\)
and value zero at \(\gamma_i\).
It is finite, independently of \(N\), by
Proposition~\ref{prop:selmer-coordinates}.
Evaluating at the two tame generators for every
prime in \(J\) shows
\[
 \length_A X_{i,\pi a}(J)
       \le\length_A T_i+
                2n_iJ_1\,\length_A(A/\pi a).
\]
Indeed the kernel has zero tame values, hence is
unramified at all these new places and lies in \(T_i\).
Set
\[
 C=J_1+\sum_i
       \bigl(\length_A T_i+
              2n_iJ_1\,\length_A(A/\pi a)\bigr).
\]
Consider pencils with the same helper, with all
lengths at modulus \(a\) at most \(\beta_i\), and
with
\[
 |J|+\sum_i\length_A X_{i,\pi a}(J)\le C .
\]
They form a nonempty class. Choose one minimizing the
nonnegative integer sum in this inequality.

Fix a row \(i\) whose length at \(\pi a\) exceeds
\(\beta_i\). If all values of every cocycle are
divisible by \(\pi\), use the injection
\(A/a\to A/\pi a,\ b\mapsto\pi b\).
Pulling back gives a cocycle obeying exactly the
same ramification and relations at modulus \(a\).
Thus its length would be at most \(\beta_i\), a
contradiction. There is a cocycle with a unit
coordinate. Theorem~\ref{thm:unseen} then implies
that \eqref{eq:unseen-inclusion} fails at some
admissible new datum \(P,\tau_P,\phi_P\).

For \(b=a,\pi a\), let \(Y_b\) be the module relaxed
at \(P,\bar P\), with the old ties, and use the maps
\[
 \Phi_0z=(z(\tau_P),z(\tau_P^-)),\qquad
 \Phi_1z=(z(\phi_P),z(\phi_P^-)).
\]
Apply Lemma~\ref{lem:tilting} at both moduli.
Together their exceptional lists have at most \(4n_i\)
residues. Avoid these and zero; the residue-field
hypothesis leaves a unit \(s\) available.
Give the new prime slope \(s^{-1}\) on row \(i\)
and zero on the other rows.

For row \(i\) its new tied module is
\(\ker(s\Phi_0-\Phi_1)\cap Y_b\). The tilting bound
puts its length below that of
\[
 Y_b\cap\ker\Phi_0\cap\Phi_1^{-1}(\Phi_0Y_b).
\]
Zero tame values identify \(Y_b\cap\ker\Phi_0\)
with the old module. At \(b=\pi a\) the last
submodule is proper, exactly by the failure of
\eqref{eq:unseen-inclusion}; its length drops by
at least one. At \(b=a\) it is a submodule of the
old module. For every other row, slope zero
requires zero tame values, so its modules can
only shrink at both moduli.
The lower-modulus constraints are preserved, the
sum at \(\pi a\) decreases by at least one, and the
one added prime leaves the potential at most \(C\).
This contradicts minimality.
All rows therefore satisfy the desired bounds.
The uniform choice \(J_0=C\) completes the induction.
\end{proof}
\src{New/Selmer/TwistTheorem.lean}{exists_isPencil}


\section{From pencils to one compatible twist}\label{sec:pencil-twist}
We prove the specialization construction with its quantifiers
and its two different sorts of points.
For finitely many balanced rows \(\rho_i\), choose an element
\(z\) in the normal closure of inertia at \(S_0\) away from
\(\ell\). Suppose the finite numbers \(B_i,c\) satisfy
\begin{equation}\label{eq:twist-anchor}
 \frac{\beta_i+
    v_A\det(\rho_i(z)-1)+v_A\det(\rho_i(z^-)-1)}
             {v_A(\ell)}\le B_i .
\end{equation}
We also allow finitely many value families \(V_k(I,\psi)\).
Their hypotheses are exactly these:
\begin{enumerate}[label=(\roman*)]
\item there are finitely many continuous characters
\(\mu_{k,t}:G\to\overline{\Q}_\ell^\times\);
\item for each auxiliary \(I\), the value is a finite sum
\(\sum_\alpha b_\alpha\psi(g_\alpha)\), with integral
coefficients \(b_\alpha\);
\item \(V_k(\varnothing,1)\ne0\), and adding an auxiliary
place \(P\) increases its valuation by at most \(8d\)
if
\(v_\ell(1-\mu_{k,t}(\phi_P)\psi(\phi_P)^\epsilon)\le d\)
for every \(t\) and \(\epsilon=\pm1\);
\item \(v_\ell(\mu_{k,t}(z)-1)\le c\).
\end{enumerate}
All coefficients and character values needed in any finite
construction lie in a sufficiently large finite extension of
\(\Q_\ell\). A continuous character on the compact group \(G\)
has unit absolute value: its absolute values form a compact
subgroup of the positive reals, and that subgroup is trivial.

\begin{lemma}[points on a torsion line]\label{lem:torsion-line}
If \(u^{\ell^k}=1\), \(t\in A\), and \(y=1+t(u-1)\), then
\[
 \ell\mid(y-1)^{\ell^k},\qquad
 \ell^D\mid y^{\ell^N}-1\quad(N\ge k+D).
\]
In particular \(N\ge k+2M\) makes \(y\) an
\(\ell^N\)-th root of unity modulo \(\ell^{2M}\),
and \(y^{\ell^{N-M}}\equiv1\pmod{\ell^M}\).
\end{lemma}
\begin{proof}
In characteristic \(\ell\), the binomial theorem gives
\((1+X)^{\ell^k}=1+X^{\ell^k}\).
Since \(u^{\ell^k}=1\), it follows that
\((u-1)^{\ell^k}\equiv0\pmod\ell\).
Multiplication by \(t^{\ell^k}\) proves the first assertion.
The same identity gives \(y^{\ell^k}\equiv1\pmod\ell\).
If \(x\equiv1\pmod{\ell^d}\), \(d\ge1\), expansion of
\(x^\ell=(1+\ell^d b)^\ell\) shows
\(x^\ell\equiv1\pmod{\ell^{d+1}}\).
For the terms with \(1<j<\ell\) use
\(\ell\mid\binom\ell j\), and for the last term use
\(d\ell\ge d+1\). This includes \(\ell=2\).
Iteration proves the claimed divisibilities.
\end{proof}
\src{New/Selmer/LinePointPowers.lean}{dvd_one_add_mul_pow_sub_one}

\begin{lemma}[the price of removing ties]\label{lem:remove-ties}
Fix a pencil and row \(i\). Let \(u^{\ell^k}=1\),
\(N\ge k+2M\), and put
\[
 y_{P_0}=u,\qquad
 y_v=1+t_{v,i}(u-1)\quad(v\ne P_0).
\]
Let \(r_y\) be the matrix function obtained by these
point coordinates. Removing the ties at a subset
\(I\subseteq J\setminus\{P_0\}\) raises its Selmer
length modulo \(\ell^M\) by at most
\[
 \sum_{v\in I,\ \epsilon=\pm}
             v_A\det(u\rho_i(\phi_v^\epsilon)-1).
\]
At \(u=1\) one may also remove the helper tie;
its cost is bounded by the two determinants for
\(\rho_i(\phi_{P_0}^\pm)-1\).
The assertions hold after any extension of coefficient DVRs.
\end{lemma}
\begin{proof}
By Lemma~\ref{lem:torsion-line}, \(r_y\bmod\ell^M\)
is a genuine representation, though \(r_y\) itself
need not be one. At a new prime the linear inertia
action is multiplication by \(y_v\).
A cocycle \(b\) defines the affine action
\(w\mapsto b(g)+r_y(g)w\) on \((A/\ell^M)^{n_i}\).
Its tame generator has order dividing \(\ell^N\).
Indeed its linear part has order dividing
\(\ell^{N-M}\) by the preceding lemma; taking that
power leaves a translation, whose \(\ell^M\)-th
power is the identity.
Wild inertia has trivial image, since it is a
pro-\(p\) group for \(p\ne\ell\), whereas the affine
inertia image is an \(\ell\)-group.

The relation between Frobenius and tame inertia is
conjugation by the \(NP\)-th power. Since
\(NP\equiv1\pmod{\ell^N}\), their affine images commute.
Evaluation at zero gives
\[
 (r_y(\phi_v^\epsilon)-1)b(\tau_v^\epsilon)
   =(r_y(\tau_v^\epsilon)-1)b(\phi_v^\epsilon).
\]
For a nonhelper place, the pencil identities give
\(r_y(\phi_v^\epsilon)=u\rho_i(\phi_v^\epsilon)\)
and \(r_y(\tau_v^\epsilon)=y_v\).
If \(t_{v,i}\ne0\), the original Frobenius acts
as the identity modulo \(\ell^M\); substituting
\(y_v-1=t_{v,i}(u-1)\) therefore shows
\[
 (u\rho_i(\phi_v^\epsilon)-1)
    \bigl(b(\tau_v^\epsilon)
                   -t_{v,i}b(\phi_v^\epsilon)\bigr)=0.
\]
If the slope is zero, the same assertion follows
directly from the commutation identity.
At the helper with \(u=1\) the identical argument
works: either the slope is zero or its original
Frobenius is trivial, making the asserted matrix
kernel the whole coefficient module.

The defect maps in parentheses have simultaneous
kernel equal to the module with the removed ties
restored. The quotient by that kernel injects into
the product of the matrix kernels just displayed.
Smith normal form over \(A\) gives, for every
square matrix \(T\),
\[
 \length_A\ker(T:(A/\ell^M)^n\to(A/\ell^M)^n)
          \le v_A(\det T);
\]
if the determinant vanishes the right side is
infinite. Add the kernel bounds. The proof uses
only divisibility and finite-length exact
sequences, so is unchanged over a larger DVR.
\end{proof}
\src{New/Selmer/TiltingLemma.lean}{IsPencil.length_le_of_drop_ties}

\begin{theorem}[one anticyclotomic character for all powers]
\label{thm:abstract-compatible-twist}
Under the displayed balanced-row and value-family
hypotheses, for every integer \(m\ge1\) and every
\(\eta>0\) there are an auxiliary set \(J\) and a
finite-order anticyclotomic character \(\psi\),
trivial at the places above \(\ell\), such that
for \(1\le d\le m\),
\begin{align}
 v_\ell V_k(J,\psi^d)
      &\le v_\ell V_k(\varnothing,1)+8c+\eta,
                              \label{eq:compatible-value}\\
 \sum_{v\in J}v_\ell(\psi^d(\tau_v)-1)
      &\le\eta,               \label{eq:compatible-conductor}\\
 \lambda_\ell\mathcal S\bigl(\psi^d\rho_i,
 S\cup J\cup\bar J,
 \mathcal L_i\cup\{b(\gamma_i)=0\};A'/\ell^M\bigr)
      &\le B_i+\eta .         \label{eq:compatible-selmer}
\end{align}
The last assertion holds over every sufficiently large
coefficient ring \(A'\), with the evident scalar
extension of the representation and relations.
The same \(J,\psi\) work for every \(i,k,d\).
\end{theorem}
\begin{proof}
Shrink a neighborhood \(W\) of \(z\) so that
\(\rho_i(\phi^\pm)\bmod\ell^M\) are constant there,
and \(v_\ell(\mu_{k,t}(\phi)-1)\le c\) there.
The latter is an open condition on \(\phi\), since
it is a lower bound for a nonzero absolute value.
Let \(J_0\) be the uniform bound of
Theorem~\ref{thm:bounded-pencil} at modulus \(\ell^M\).
Choose \(\delta>0\) with
\begin{equation}\label{eq:delta-twist}
 (8J_0+2)\delta<\eta.
\end{equation}
All uses of the torsion-point theorem below are
uniform in polynomial degree and coefficient ring.

First choose \(k\) large enough for one-variable
simultaneous evaluation with one power, the
finitely many bounds \(B_i\) and
\(v_\ell V_k(\varnothing,1)+8c\), tolerance
\(\delta\), and the following upper bound on the
total degree of the monic avoidance polynomials:
\[
 2J_0\sum_i n_i+2J_0\sum_k|T_k|.
\]
Then choose \(N\ge k+2M\), large enough for
Theorem~\ref{thm:many-torsion} with at most \(J_0\)
variables, powers \(1,\ldots,m\), tolerance \(\delta\),
the finite bounds just mentioned plus \(8J_0\delta+\delta\),
and monic degree sum \(m(m+1)/2\).
Here the closeness exponent in that theorem is
\(\ell^k\), fixed before \(N\).
Choose a pencil at this \(N\) with \(|J|\le J_0\).

Write \(\theta_v=h_v\).
For any tuple of actual \(\ell^N\)-th roots define
\(\psi_y(g)=\prod_v y_v^{\widetilde\theta_v(g)}\).
It is a character, anticyclotomic and induced by
\(x/\bar x\) on ideles. If the nontrivial coordinates
are in \(I\), it has the required modulus \(I\).
The value-family formula supplies integral
polynomials \(V_{k,I}(y)\) for all \(I\subseteq J\),
after one finite coefficient extension. Explicitly
replace each \(\psi_y(g_\alpha)\) in (ii) by
\(\prod_v y_v^{\widetilde\theta_v(g_\alpha)}\).
The finite number of subsets causes no uniformity
problem.

For each row choose a Fitting determinant polynomial
\(G_i(y)\) for the Selmer module with ties at
\(J\setminus\{P_0\}\), sharp at \(y=1\).
The finite-coordinate proposition justifies its
existence. The pencil bounds the module with all
ties by \(\beta_i\). Removing the helper tie at
\(y=1\), and using the constancy of its Frobenius
modulo \(\ell^M\), bounds the remaining module
by the length appearing in
\eqref{eq:twist-anchor}.
Consequently \(v_\ell G_i(1)\le B_i\).

Restrict \(G_i\) to the row's own line
\[
 y_{P_0}=X,\qquad y_v=1+t_{v,i}(X-1).
\]
For value family \(k\), use instead the line
\(y_{P_0}=X,\ y_v=1\) for \(v\ne P_0\).
At \(X=1\), its helper-only polynomial
\(V_{k,\{P_0\}}\) has valuation at most
\(v_\ell V_k(\varnothing,1)+8c\), by (iii)
at the helper. Apply the one-variable
torsion theorem simultaneously to these line
polynomials and to the monic polynomials
\[
 \det(X-\rho_i(\phi_v^\epsilon)^{-1}),\qquad
 X-\mu_{k,t}(\phi_v)^\epsilon .
\]
The diagonal scalar notation in the determinant
means \(XI-\rho_i(\phi_v^\epsilon)^{-1}\).
We obtain a common \(u^{\ell^k}=1\), with losses
at most \(\delta\) in all the line polynomials
and valuation at most \(\delta\) for every
avoidance polynomial. Coefficients and \(u\)
may be put into one finite coefficient ring.

Let \(y^{(i)}\) be the point on row \(i\)'s line
at \(u\). It is generally only a torsion point
modulo \(\ell^M\); we use the matrix-function
definition of \(\mathcal S\) at this point.
Lemma~\ref{lem:remove-ties} removes the remaining
ties at cost at most \(2J_0\delta\).
The determinant estimates follow from the monic
characteristic polynomials above, since
\(\det\rho_i(\phi_v^\epsilon)\) is a unit.
Thus the normalized untied length at \(y^{(i)}\)
is at most
\(B_i+(2J_0+1)\delta\).
Now choose an untied Fitting polynomial \(F_i\)
\emph{sharp at this particular point}
\(y^{(i)}\). Its value there satisfies the same
bound, and its values everywhere bound the
untied Selmer lengths from above. The sharpness
point must be chosen after \(u\); sharpness at
\(1\) would give the wrong direction for this step.

For a value family put \(y^{(k)}_{P_0}=u\) and
\(y^{(k)}_v=1\) elsewhere. This tuple is an actual
torsion tuple. At a nonhelper Frobenius its
character has value \(u\).
The monic estimates imply
\[
 v_\ell(1-\mu_{k,t}(\phi_v)u^\epsilon)\le\delta.
\]
For \(\epsilon=1\) use
\(X-\mu_{k,t}(\phi_v)^{-1}\), and for
\(\epsilon=-1\) use \(X-\mu_{k,t}(\phi_v)\);
multiplication by the units \(u,\mu_{k,t}\)
does not affect valuation. Add the nonhelper
places one at a time using (iii). This yields
\[
 v_\ell V_{k,J}(y^{(k)})
 \le v_\ell V_k(\varnothing,1)+8c+
                         (8J_0+1)\delta.
\]
Every anchor coordinate, of either type, satisfies
\(\ell\mid(y_v-1)^{\ell^k}\).

Apply Theorem~\ref{thm:many-torsion} with the
polynomials \(F_i,V_{k,J}\), using the appropriate
anchor \(y^{(i)}\) or \(y^{(k)}\) for every power
\(d\). The polynomials are the same for each
\(d\); their anchors are allowed to be different
for different indices. Include
\(D_d(X)=X^d-1,\ 1\le d\le m\), among the
monic avoidance polynomials.
We obtain a single tuple \(\zeta\) with
\[
 \zeta_v^{\ell^N}=1,\qquad
 v_\ell(\zeta_v^d-1)\le\delta\quad(1\le d\le m),
\]
and all polynomial valuations at \(\zeta^d\)
bounded by their anchor bounds plus \(\delta\).
Set \(\psi=\psi_\zeta\).
The pencil gives
\(\psi^d(\tau_v)=\zeta_v^d\);
hence the conductor sum is at most
\(J_0\delta<\eta\).
In particular none of these tame values is
\(1\), so no power loses a required conductor
place.

The value loss is at most
\((8J_0+2)\delta<\eta\), proving
\eqref{eq:compatible-value}.
For each row the untied Fitting polynomial
bounds the length at the actual character
\(\psi^d\) by
\(B_i+(2J_0+2)\delta\le B_i+\eta\).
The determinant bound holds under every
further extension of the coefficient DVR,
and normalized valuations are unchanged.
This proves \eqref{eq:compatible-selmer}
with its eventual-ring quantifier.
\end{proof}
\src{New/Selmer/TwistTheorem.lean}{exists_anticyclotomic_twist}

The special value hypotheses, the balancing of the particular
automorphic rows, and the constants in
Input~\ref{in:twist} still have to be verified from their
definitions. The theorem just proved supplies the complete
global selection argument once those local and value inputs
are given. It does not assume that a Fitting polynomial at a
general tuple defines a characteristic-zero Galois character.

\section{Local losses and their global budget}\label{sec:budget}
Fix admissible data \(D\). Write \(f\) for the number of infinite places of
its CM field \(E\), and retain the source's positive integers
\(a,m,k,t,h_2,b_1\), the auxiliary product \(g=g_0g_1\), and scale \(b>1\).
The normalizations are
\[
 \mathcal B=f a(\log b)^3,\qquad B_0=\mathcal B/(mk).
\]
In the source \(k=r+1-m\), \(t=5r\), \(h_2=r^8\), \(a\) is the least multiple
of \(N_1\) above \(xb\), and \(b_1\) the least multiple of \(N_1\) above \(a/x\).

Let
\[
 G=\prod_{w\ {\rm chosen}}\#\mathrm{GL}_{rt}(\F_w).
\]
There are \(f\) chosen places above each \(g_i\), all with residue field of
cardinality \(g_i\). This integer pays for stabilizers, finite idempotents,
and the scalar inertia elements used in the denominator argument.
Put
\begin{align}
 J&=(2b_1)!\left(\prod_{\max(b_1-30h_2^3,0)<s\le b_1+30h_2^3}s\right)^{h_2^3},
                                                     \label{eq:J}\\
 c_\ell&=v_\ell(G)+2^{-\ell},\\
 \lambda_\ell&=10^6\frac m k\left(b_1v_\ell(g)+v_\ell(J)\right),\\
 \beta_\ell&=h_2^{27}\left(c_\ell+v_\ell(X)\right)+4f\lambda_\ell .
                                                     \label{eq:beta}
\end{align}
The integer \(X\ge1\), independent of \(\ell\), is supplied by the odd
Tate--Shafarevich estimate for the abelian CM field \(E\). With
\(N=N_1d\) its cyclotomic conductor, its size bound is
\begin{equation}\label{eq:X}
 X\le 2^{4(N+t+3)^{27}}h_E ,
\end{equation}
where \(h_E\) is the class number. The source theorem controls the specified
odd Tate-twisted cohomology modules simultaneously for every prime,
including \(2\), every \(0\le e<t\), and every modulus \(\ell^\nu\).
A full definition and proof of that cohomological assertion are still owed
in this manuscript; here only its uniform positive integer and size bound
are used.
\src{New/Selmer/DeterminantBound.lean}{lambdaLoss; loss}
\src{New/Selmer/AptTupleDefs.lean}{crudeUnit}
\src{New/Selmer/RibetElement.lean}{rowBudget}
\src{Lib/Cyclotomic/OddSha/OddShaBound.lean}{exists_length_shaOdd_le}

\begin{lemma}\label{lem:prime-log}
For a positive integer \(A\) and any finite set \(T\) of primes,
\[
 \sum_{\ell\in T}v_\ell(A)\log\ell\le\log A .
\]
Also \(\sum_{\ell\in T}2^{-\ell}\log\ell\le2\).
\end{lemma}
\begin{proof}
The full prime factorization of \(A\), after taking logarithms, gives
\(\log A=\sum_\ell v_\ell(A)\log\ell\); every summand is nonnegative.
For the second inequality use \(\log\ell\le\ell\) and
\(\sum_{n\ge0}n2^{-n}=2\).
The latter identity follows by multiplying
\(\sum_{n\ge0}n z^n\) by \(1-z\): subtracting the shifted sum leaves
\(\sum_{n\ge1}z^n=z/(1-z)\), hence the original sum is \(z/(1-z)^2\).
The power series converge absolutely for \(\abs z<1\).
\end{proof}

\begin{proposition}[the prime sum is finite without choosing one twist globally]
\label{prop:budget}
For every finite set \(T\) of primes,
\begin{align}
 \sum_{\ell\in T}c_\ell\log\ell&\le\log G+2,             \label{eq:cr-sum}\\
 \sum_{\ell\in T}\beta_\ell\log\ell
 &\le h_2^{27}(\log G+\log X+2)
       +4\cdot10^6f\,\frac m k(b_1\log g+\log J).       \label{eq:budget-sum}
\end{align}
Moreover
\[
 \log J\le
 2b_1\log(2b_1)+60h_2^6\log(b_1+30h_2^3),
 \qquad
 \log G\le (rt)^2 f\log g .
\]
\end{proposition}
\begin{proof}
Insert \eqref{eq:beta}, distribute the finite sum, and apply
Lemma~\ref{lem:prime-log} separately to \(G,X,g,J\).
This proves \eqref{eq:cr-sum} and \eqref{eq:budget-sum}, with no interchange
of an infinite sum.

There are at most \(60h_2^3\) integers in the interval of \eqref{eq:J}.
Each is positive and at most \(b_1+30h_2^3\).
Also \((2b_1)!\le(2b_1)^{2b_1}\). Taking logarithms proves the first bound.
For the second, a matrix group of size \(s\) over \(\F_q\) has at most
\(q^{s^2}\) elements, because it is a subset of all \(s\)-by-\(s\) matrices.
Multiply over the chosen places and use
\(\prod_w\#\F_w=g^f\).
\end{proof}

The global arithmetic bookkeeping is therefore reduced to explicit bounds
for logarithms of positive integers. For example the following elementary
calculation isolates the potentially large factorial term.

\begin{lemma}[absorption of the local factorial loss]\label{lem:factorial-budget}
Assume
\[
 \begin{gathered}
 b\ge4,\quad \log b\ge10^6,\quad b_1\le2b,\quad 60h_2^6\le b,\\
 30h_2^3\le b,\quad \log g\le22\log b,\quad m^2\le4x,\quad a\ge xb .
 \end{gathered}
\]
Then
\[
 4\cdot10^6 f\,\frac m k(b_1\log g+\log J)\le \frac12 B_0 .
\]
\end{lemma}
\begin{proof}
The bounds for \(J\) in Proposition~\ref{prop:budget} give
\[
 b_1\log g+\log J
 \le 44b\log b+4b\log(4b)+b\log(3b)\le54b\log b.
\]
Here \(\log(4b)\le2\log b\) and \(\log(3b)\le2\log b\) since \(b\ge4\).
The ratio of the resulting bound to \(B_0=f a(\log b)^3/(mk)\) is at most
\[
 \frac{216\cdot10^6 m^2b}{a(\log b)^2}
 \le \frac{864\cdot10^6}{(\log b)^2}<\frac12 .
\]
\end{proof}

Section~\ref{sec:numerics} proves the coarse inequalities from admissibility,
bounds \(h_2^{27}(\log G+\log X+2)\) by \(B_0/2\), and proves
\(9(\log G+2)\le x^{999}\). This closes both prime budgets, conditional
only on the stated size bound for the integer \(X\).
The cohomological theorem producing \(X\) remains separate.
\src{New/Selmer/FitTwistsExist.lean}{the two finite-prime budget bounds in exists_isFitFamily}



\section{Closing the numerical budgets}\label{sec:numerics}
We now discharge the scale comparisons used above. The odd-class theorem
still supplies \(X\) and \eqref{eq:X}; the argument here proves that its
stated size, together with the local losses, fits the required budget.

\subsection{A crude class-number bound}
We use the classical Minkowski ideal-class theorem in its precise usual
form: in a degree-\(n\) number field of discriminant \(\Delta\) and
\(r_2\) complex pairs, every ideal class has an integral representative
of norm at most
\((4/\pi)^{r_2}(n!/n^n)\sqrt{\abs\Delta}\).
This is an established geometry-of-numbers input; no assertion about
RH is used.

\begin{lemma}\label{lem:class-number}
For a number field \(F\) of degree \(n\ge1\),
\[
 h_F\le (2^n\abs{\Delta_F})^{n+1}.
\]
\end{lemma}
\begin{proof}
Put \(B=2^n\abs{\Delta_F}\), a positive integer.
Minkowski's bound is at most \(B\): use \(\pi\ge2\), \(r_2\le n\),
\(n!\le n^n\), and \(\sqrt{\abs\Delta}\le\abs\Delta\).
Thus every ideal class has an integral ideal \(J\) of positive norm
\(a\le B\).

Every such ideal contains \(a\mathcal O_F\), since the finite additive
group \(\mathcal O_F/J\) has order \(a\) and is killed by \(a\).
It has the form \(J=(a,z)\).
Here are details of this last fact.
For each prime ideal \(\mathfrak p\) dividing \(a\mathcal O_F\), choose
\(z_{\mathfrak p}\in J\setminus\mathfrak pJ\);
this set is nonempty, since otherwise multiplying
\(J=\mathfrak pJ\) by the inverse ideal would give
\(\mathcal O_F=\mathfrak p\).
The Chinese remainder theorem for the finitely many distinct prime
ideals gives an element \(z\in J\) congruent to \(z_{\mathfrak p}\)
modulo \(\mathfrak pJ\) for each of them: use the ordinary scalar
Chinese-remainder coefficients and sum their multiples of the
\(z_{\mathfrak p}\).
At these primes \(v_{\mathfrak p}(z)=v_{\mathfrak p}(J)\).
At primes not dividing \(a\), both \(J\) and \((a,z)\) have valuation zero.
Unique ideal factorization therefore gives \(J=(a,z)\).

For fixed \(a\), the ideal \((a,z)\) depends only on \(z\) modulo
\(a\mathcal O_F\), which has \(a^n\) elements, since \(\mathcal O_F\)
is free of rank \(n\) over \(\Z\).
Consequently the number of ideals needed to represent the classes is at most
\(\sum_{a=1}^B a^n\le B^{n+1}\).
\end{proof}
\src{Lib/General/NumberTheory/ClassNumberBound.lean}{classNumber_le}

\subsection{Coarse consequences of admissibility}
Write
\[
 U=\log x,\qquad V_1=\log x_1,\qquad B=\log b.
\]
The data have \(c\ge64\), \(N_0,d\le c\), \(s\ge1\), and
\begin{equation}\label{eq:scales-explicit}
 \begin{gathered}
 10^{12}c\le\log B,\quad 100\log B\le U,\quad
 10^6U\le V_1,\quad 10^6V_1\le B,\quad s\le32U,\\
 \sqrt x\le m\le2\sqrt x,\qquad x^5\le r\le3x^5 .
 \end{gathered}
\end{equation}
The remaining definitions are
\[
 k=r+1-m,\quad t=5r,\quad h_2=r^8,\quad
 N_1=N_0\,3^{\lfloor\log_3s\rfloor+5}2^s,\quad N=dN_1 .
\]
Rounding gives
\(xb\le a\le xb+N_1\) and
\(a/x\le b_1\le a/x+N_1\).
The two level primes lie in \([b^{10},b^{11}]\).

\begin{proposition}\label{prop:coarse}
For these data,
\[
 \begin{gathered}
 U\ge100,\quad B\ge10^6,\quad
 B\ge10^{12}U,\quad B\le x^{1/100},\quad b\ge x^{10^{12}},\\
 N_1\le N\le x^{24},\quad 1\le f\le x^{24},\quad
 r/2\le k\le r,\quad m\le x,\quad r,t\le x^6,\quad h_2\le x^{41},\\
 b_1\le2b,\quad 60h_2^6\le b,\quad 30h_2^3\le b,\quad
 \log g\le22B,\quad
 \operatorname{rd}(E)\le N\le x^{24}.
 \end{gathered}
\]
\end{proposition}
\begin{proof}
The first two scale inequalities imply
\(U\ge100\cdot10^{12}c\ge100\) and \(B\ge e^{10^{12}c}>10^6\).
The next two give \(B\ge10^{12}U\); exponentiating gives
\(b\ge x^{10^{12}}\).
Also \(100\log B\le U\) gives \(B\le x^{1/100}\).

Since \(3^{\lfloor\log_3s\rfloor+5}\le243s\),
\[
 \log N\le2\log c+\log(7776U)+32U\log2.
\]
Here \(2\log c\le2\log B\le U/50\),
\(\log(7776U)\le U\) for \(U\ge100\), and \(\log2<7/10\).
The last two elementary inequalities follow from monotonicity of
\(U-\log U\), and the exponential series, respectively.
Thus \(\log N\le(1/50+1+112/5)U<24U\).
In particular \(N_1\le x^{24}\).
As \(E=K(\zeta_{N_1})\), with \([K:\Q]=2\), its degree is at most
\(2\varphi(N_1)\). It is totally imaginary, so
\(f=[E:\Q]/2\le\varphi(N_1)\le N_1\).

For \(x\ge e^{100}\), the rank bounds give \(m\le r/2\), whence
\(r/2\le k\le r\); they also give
\(m\le x\), \(r,t\le x^6\), and
\(h_2\le3^8x^{40}\le x^{41}\).
Rounding gives
\(b_1\le b+N_1/x+N_1\le2b\), since
\(N_1\le x^{24}\le b/4\) and \(x\ge1\).
Further
\(60h_2^6\le60x^{246}\le x^{247}\le b\) and
\(30h_2^3\le30x^{123}\le x^{124}\le b\).
The bounds on \(g_0,g_1\) give \(\log g\le22B\).

For the discriminant bound, \(K=\Q(\sqrt{-d})\) lies in
\(\Q(\zeta_d)\). This inclusion has a direct Gauss-sum proof:
for the quadratic character \(\chi\) modulo the prime \(d\), put
\(\tau=\sum_{a\ne0}\chi(a)\zeta_d^a\).
One has \(\bar\tau=\chi(-1)\tau=-\tau\), since \(d\equiv7\pmod8\).
Also
\[
 \tau\bar\tau
 =\sum_{u\ne0}\chi(u)\sum_{b\ne0}\zeta_d^{(u-1)b}
 =(d-1)-\sum_{u\ne0,1}\chi(u)=d.
\]
Thus \(\tau^2=-d\).
It follows that \(E\subseteq\Q(\zeta_{dN_1})\).
The classical cyclotomic discriminant formula is
\[
 \abs{\Delta_{\Q(\zeta_N)}}
 =N^{\varphi(N)}
    \prod_{p\mid N}p^{-\varphi(N)/(p-1)}
 \le N^{\varphi(N)}.
\]
For a subfield \(F\subseteq L\), the relative discriminant identity
\(\abs{\Delta_L}=\abs{\Delta_F}^{[L:F]}
                 N_{F/\Q}(\mathfrak d_{L/F})\)
and integrality of the relative discriminant give
\(\operatorname{rd}(F)\le\operatorname{rd}(L)\).
Apply this with \(F=E,L=\Q(\zeta_N)\).
\end{proof}
\src{New/Data/AdmissibleData.lean}{Data; Data.k; Data.roundUp}
\src{New/Data/Different.lean}{towerDeg; Parameters.E}
\src{New/Data/Magnitudes.lean}{the admissible scale comparisons}

The proposition uses weaker coarse bounds than the source's sharper
\(f\) and root-discriminant bounds.

\begin{theorem}[the two global budgets]\label{thm:global-budgets}
Assume the positive integer \(X\) satisfies \eqref{eq:X}.
For every finite set \(T\) of primes,
\[
 \sum_{\ell\in T}\beta_\ell\log\ell\le B_0,\qquad
 \sum_{\ell\in T}9c_\ell\log\ell\le x^{999}.
\]
\end{theorem}
\begin{proof}
We bound the first term of \eqref{eq:budget-sum}.
By Proposition~\ref{prop:coarse},
\[
 \log G\le(rt)^2 f\,22B\le x^{24}x^{24}x^2=x^{50}.
\]
The last factor uses \(22B\le x^2\), valid for \(x\ge e^{100}\).

Put \(n=[E:\Q]=2f\). Both \(n,n+1\le x^{25}\).
Lemma~\ref{lem:class-number} and \(\operatorname{rd}(E)\le x^{24}\) give
\[
 \log h_E\le n(n+1)(\log2+\log\operatorname{rd}(E))
            \le x^{50}(\log2+24U)\le x^{52}.
\]
Also \(N+t+3\le3x^{24}\le x^{25}\).
Hence \eqref{eq:X} gives
\[
 \log X\le4x^{675}\log2+x^{52}\le x^{676}.
\]
As \(h_2^{27}\le x^{1107}\), we get
\[
 h_2^{27}(\log G+\log X+2)
       \le x^{1107}x^{677}=x^{1784}.
\]
On the other hand \(mk\le x^7\), \(a\ge xb\), \(f\ge1\), and \(B\ge1\),
so
\[
 B_0=\frac{faB^3}{mk}\ge f b x^{-6}
       \ge x^{10^{12}-6}\ge2x^{1784}.
\]
Thus this part costs at most \(B_0/2\).
All hypotheses of Lemma~\ref{lem:factorial-budget} now follow from
Proposition~\ref{prop:coarse} and \(m^2\le4x\).
That lemma bounds the other part of \eqref{eq:budget-sum} by \(B_0/2\).
Adding proves the first assertion.

Finally, by \eqref{eq:cr-sum},
\[
 \sum_{\ell\in T}9c_\ell\log\ell
       \le9(\log G+2)\le9(x^{50}+2)\le x^{52}\le x^{999}.
\]
All estimates hold for every finite prime set, before a twist is chosen.
\end{proof}
\src{New/Selmer/FitTwistsExist.lean}{the finite-prime budget estimates}

\begin{proposition}[the scale comparisons in Theorem C]\label{prop:C-scales}
With \(\mathcal B=faB^3\), \(B_0=\mathcal B/(mk)\),
and \(L=B\mathcal B\), admissibility implies
\[
 b\le L,\qquad
 2x_1^2b^2+4m^2t x^{999}\le b^3/4,\qquad
 3rB_0+10m x^{999}+6kB_0\le k B B_0.
\]
\end{proposition}
\begin{proof}
Since \(f\ge1\), \(a\ge xb\ge b\), and \(B\ge1\), one has \(L=faB^4\ge b\).
Further \(B\ge10^6V_1\), \(V_1\ge1\), imply
\(B-2V_1\ge\log16\). Exponentiating yields
\(2x_1^2b^2\le b^3/8\).
For the other term, \(m\le x\), \(t\le x^6\) give
\(4m^2t x^{999}\le4x^{1007}\).
The inequalities
\(\log32+1007U\le3B\), from \(B\ge10^{12}U\) and \(U\ge100\),
give \(4x^{1007}\le b^3/8\). Add the two bounds.

Finally \(r\le2k\) gives \(3rB_0+6kB_0\le12kB_0\), while
\[
 \frac{10m x^{999}}{kB_0}
    =\frac{10m^2x^{999}}{faB^3}
    \le\frac{40x^{999}}{bB^3}\le1
\]
by \(m^2\le4x\), \(a\ge xb\), \(f\ge1\), and \(b\ge x^{10^{12}}\).
Thus the sum is at most \(13kB_0\le k B B_0\).
\end{proof}
\src{New/Endgame/LongForms.lean}{the scale comparisons used in Theorem C}

\section{The compatible family and the denominator input}\label{sec:fit}
For a coefficient ring \(\mathcal O\) finite over \(\Z_\ell\), write
\[
 \lambda_\ell(M)=\length_{\mathcal O}(M)/v_{\mathcal O}(\ell)
\]
for normalized length. This removes dependence on the ramification
degree of the coefficient ring.

A continuous \(1\)-cocycle for \(\rho:G_E\to\mathrm{GL}(L)\) is a function
\(z:G_E\to L/\ell^M L\) with \(z(gh)=z(g)+\rho(g)z(h)\).
The Selmer cocycles used here are unramified off
\[
 \{\text{places above }\ell\}\ \cup\
 \{\text{chosen places above }g\}\ \cup Q\cup\bar Q,
\]
and obey the linear relations satisfied by crystalline cocycles at
the places above \(\ell\). One also imposes \(z(\gamma)=0\) at a specified
scalar element. Thus it is a module of actual cocycles with stated
relations, not an unnamed quotient of a cohomology group.

The admissible Galois rows have prescribed jumps, conjugate self-duality,
nondegenerate trace, scalar inertia, and full matrix span on sufficiently
large coefficient lattices. These hypotheses are essential to the
reciprocity argument which makes the Selmer groups short.

\begin{sourceinput}[one twist for finitely many rows]\label{in:twist}
Fix \(D,\ell,\iota\), finitely many admissible rows \(R_i\) of ranks \(r_i\),
and an integer \(M\ge0\). There is one finite-order weight-zero tame
character \(\tau\) whose first \(m\) powers are allowed twists. For
\(1\le j\le m\), put \(\vartheta=\tau^j\), and let \(\psi\) be the
\(\ell\)-adic avatar of \(\vartheta^{-1}\).
\begin{enumerate}[label=(\roman*)]
\item Its set \(Q\) of new ramified places has prime norms, is disjoint
from its conjugate and the prescribed bad places, and is prime to \(\ell\).
The character comes from the norm-one torus, is trivial at the required
places above \(\ell\), and satisfies
\(\sum_{w\in Q}v_\ell(\psi(\sigma_w)-1)\le c_\ell\)
for the chosen tame inertia generators.
\item For every \(i\) there are a scalar Galois element \(\gamma_i\)
and a scalar \(u_i\) with \(u_i-1\) a unit, \(\mathrm{cyc}(\gamma_i)=1\),
and \(\psi R_i(\gamma_i)=u_i I\), such that, over all sufficiently
large coefficient rings, the Selmer-cocycle module just described,
with \(z(\gamma_i)=0\), has normalized length at most
\((10f r_i+1)c_\ell\).
\item For every ramified pair \(p\) of the data and each sign \(s\),
the source's ratio \(\rho_p(\vartheta,s)\) of twisted to untwisted Hecke
values satisfies
\(\abs{\iota(\rho_p(\vartheta,s)^{-1})}_\ell
 \le\exp(9c_\ell\log\ell)\).
\end{enumerate}
The same \(\tau\) works for all \(i,j,p,s\) in these statements.
\end{sourceinput}
\src{New/Selmer/TwistForRows.lean}{exists_mFoldTwist_selmerLengthLe}

Theorem~\ref{thm:abstract-compatible-twist} proves the global selection
argument, including finite coordinates, reciprocity, and the shortening
by bounded pencils. To discharge the present input one must still verify
the local balancing estimates and the value-family hypotheses for the
specified automorphic rows and Hecke ratios. These source-specific
verifications are not consequences of polynomial avoidance.

\begin{sourceinput}[denominator construction]\label{in:denominator}
Let \(\vartheta\) satisfy the auxiliary conditions in
Input~\ref{in:twist}(i). Let \(O\) be an auxiliary pair of ranks
\((r_O,n_O)\), \(e_0\) its specified group-algebra element, and
\(Z\subseteq\mathcal F(O)\). Suppose \(Z\) is spanned by vectors in \(Z\)
with both of the following properties:
\begin{enumerate}[label=(\roman*)]
\item the vector lies in the projection for a tuple of apt labels for
\((O,\vartheta)\), and that projection has image in \(\im e_0\);
\item it lifts some Hecke data which have an admissible row for an
auxiliary pair \(V\) of the same ranks, with
\(\mu_O=\mu_V\vartheta\), and this row satisfies the short Selmer bound
of Input~\ref{in:twist}(ii) at
\(M=\lfloor r\beta_\ell\rfloor+1\).
\end{enumerate}
Then
\[
 \operatorname{Den}_\iota(O,e_0,Z;\beta_\ell\log\ell)
\]
holds in the sense of \eqref{eq:den}.
\end{sourceinput}
\src{New/Selmer/DenominatorBound.lean}{denominatorLe_rowBudget}

The integer \(r\) in the truncation is the fixed rank of \(D\), not the
dimension of \(Z\) or necessarily \(r_O\).
A complete prose proof of Input~\ref{in:denominator} must construct the
operators of Proposition~\ref{prop:operators}, with the valuation of
\(G^3\) times their eigenvalue at most \(r_O\beta_\ell\).
It must also show they clear denominators on the tested complement.
The source uses triangular Taylor--Wiles primes, a trace ideal, a Ribet
determinant, and an integral diagonal patching operator for this purpose.
None of these conditions is discharged merely by attaching the name of
the Lean declaration to this statement.

\subsection{Assembly of a fit family}
Here is the precise assembly, separated from those arithmetic inputs.
Suppose the good Hecke eigensystems have the required admissible rows,
and the source's transfer theorem identifies their occurrence across
the \(m\) spaces \(W_j\). Also include the finitely many rows for the
rank-one pairs at all indices \(j\). The eigensystem set is finite:
commuting normal Hecke operators on a finite-dimensional space have a
simultaneous orthogonal eigenbasis, so only finitely many joint systems
can occur. The rank-one pair set is finite in the admissible data.

For each \((\ell,\iota)\) apply Input~\ref{in:twist} to this one finite
union of rows and to \(M=\lfloor r\beta_\ell\rfloor+1\).
Call the resulting \(m\)-fold twist \(\theta_{\ell,\iota}\).
There is no dependence on a later choice of an individual good eigenform.
Apply Input~\ref{in:denominator} first to the image of the \(j\)-th
good adjoint theta lift and then to every rank-one line at that index.
The source's apt-label and Hecke-transfer assertions verify its two
spanning hypotheses. This gives the denominator bound at the same
\(\theta_{\ell,\iota}^j\) in both applications.

Set \(\gamma_\ell=\beta_\ell\log\ell\).
By Theorem~\ref{thm:global-budgets}, every finite prime sum of
\(\gamma_\ell\) is at most \(B_0\).
Input~\ref{in:twist}(iii), \eqref{eq:cr-sum}, and the height calculus give
\[
 \HH{x^{999}}{\left(\rho_p(\theta_{\ell,\iota}^j,s)^{-1}\right)}
\]
for every ramified \(p,s,j\).
Coprimality of the conductor with \(\ell\) follows from (i).
These three assertions are exactly the fitness requirements:
coprimality, a summable denominator bound for every good lift and
rank-one line, and a height bound for every required ramified ratio.
\src{New/Selmer/FitTwists.lean}{Data.IsFitFamily}
\src{New/Selmer/FitTwistsExist.lean}{exists_isFitFamily}

The transfer between the \(W_j\) is a substantive assertion: their
archimedean signatures differ. It is proved in the source from rigid
labels and comparison of trace formulas, and remains a separate
translation obligation here.


\section{Heights and the product formula}\label{sec:heights}
Let
\[
 \mathcal I=\{(\ell,\iota):\ell\text{ prime},\
                \iota:\C\xrightarrow{\sim}\overline{\Q}_\ell\}.
\]
These are field isomorphisms, with no continuity requirement.
A family \(z=(z_{\ell,\iota})\) may depend on both indices.

\begin{definition}\label{def:height}
Write \(\HH h z\) if there are a complex number \(H\ne0\), a real number
\(\alpha\), and real numbers \(e_\ell\) such that
\begin{align*}
 &\abs{\sigma(H)}\le e^\alpha &&(\sigma\in\Aut(\C)),\\
 &\abs{\iota(H z_{\ell,\iota})}_\ell\le e^{e_\ell}
                                      &&((\ell,\iota)\in\mathcal I),\\
 &\alpha+\sum_{\substack{\ell<T\\\ell\ {\rm prime}}}e_\ell\le h
                                      &&\text{for all sufficiently large }T.
\end{align*}
\end{definition}
The letter \(H\) is one number for the whole family.
The \(e_\ell\) need not be nonnegative, and \(h\) may be negative.
For example, for an integer \(a\ge1\) the constant family \(1/a\) has bound
\(\HH{\log a}{1/a}\), with witness \(H=a,\alpha=\log a,e_\ell=0\).
The constant family \(a\) has bound \(\HH{-\log a}a\), with
\(H=1/a,\alpha=-\log a,e_\ell=0\).
Thus this relation is not the usual nonnegative Weil height.

\begin{lemma}[height calculus]\label{lem:height-calculus}
The relation has the following properties.
\begin{enumerate}[label=(\roman*)]
\item If \(\HH h z\), \(h\le h'\), and
\(\abs{\iota(z'_{\ell,\iota})}_\ell
 \le\abs{\iota(z_{\ell,\iota})}_\ell\) at every index, then \(\HH{h'}{z'}\).
\item If all finite prime sums of \(e_\ell\) are at most \(h\), and
\(\abs{\iota(z_{\ell,\iota})}_\ell\le e^{e_\ell}\), then \(\HH h z\).
\item If \(\HH h z\) and \(\HH k w\), then \(\HH{h+k}{zw}\).
Consequently finite products add bounds, and
\(\HH h z\) implies \(\HH{nh}{z^n}\) for \(n\in\N\).
\end{enumerate}
\end{lemma}
\begin{proof}
For (i), use the same witness and multiply the local comparison by
\(\abs{\iota(H)}_\ell\). For (ii), use \(H=1,\alpha=0\).
For (iii), multiply the witnesses \(H,K\), add their archimedean exponents,
and add their prime sequences. Multiplicativity of each absolute value
gives both local bounds. Beyond the larger of the two thresholds,
the two finite-sum inequalities can be added. Induction gives finite
products and powers; the empty product has witness \(1,0,(0)_\ell\).
\end{proof}
\src{New/Heights/Heights.lean}{HeightLe.mono; HeightLe.of_le; HeightLe.mul; HeightLe.prod; HeightLe.pow}

\subsection{Why the constant family has a nonnegative height bound}
We give the field-theoretic details to avoid an implicit algebraicity assumption.

\begin{lemma}[extension of embeddings]\label{lem:extend}
An embedding of a countable characteristic-zero subfield \(F\subseteq\C\)
into \(\overline{\Q}_\ell\) extends to an isomorphism
\(\C\simeq\overline{\Q}_\ell\).
An embedding \(F\to\C\) likewise extends to an automorphism of \(\C\).
\end{lemma}
\begin{proof}
First extend the embedding to the algebraic closure of \(F\) inside the
target, using the usual maximal-extension argument: a partial embedding
on an algebraic intermediate field extends across the minimal polynomial
of any element not yet in its domain, since that polynomial has a root
in the algebraically closed target. The algebraic closures of a countable
field are countable.

Both \(\C\) and \(\overline{\Q}_\ell\) have cardinality \(2^{\aleph_0}\).
For the latter, \(\Q_\ell\) has this cardinality from its digit expansions,
and adjoining algebraic elements does not change an infinite cardinality.
The transcendence degree of either field over a countable algebraically
closed subfield is \(2^{\aleph_0}\): for a transcendence basis of infinite
cardinality \(\kappa\), its algebraic closure has cardinality
\(\max(\aleph_0,\kappa)\); a finite or countable basis would give a countable
field. Choose a bijection between the two transcendence bases.
It extends the already given embedding to an isomorphism of their rational
function fields. Extend once more across algebraic closures by the same
maximal-extension argument. An algebraically closed field algebraic over
the image is the full image, so the resulting map is onto.
\end{proof}

\begin{lemma}\label{lem:algebraic}
If a complex number \(H\) has a uniform bound
\(\abs{\iota(H)}_2\le B\) for all \(\iota:\C\simeq\overline{\Q}_2\),
then \(H\) is algebraic over \(\Q\).
\end{lemma}
\begin{proof}
If \(H\) were transcendental, choose a transcendental \(t\in\Q_2\) with
\(\abs t_2>B\). Such a choice is possible: every nonempty open ball of
\(\Q_2\) is uncountable, whereas the elements algebraic over \(\Q\) are
countable; there are balls on which the absolute value exceeds \(B\).
The map \(\Q(H)\to\Q(t)\), \(H\mapsto t\), is a field isomorphism.
Extend it by Lemma~\ref{lem:extend} to an isomorphism of \(\C\) with
\(\overline{\Q}_2\). Its value on \(H\) contradicts the bound.
\end{proof}

\begin{theorem}[product formula for a varying family]\label{thm:height-nonneg}
If \(\HH h1\), then \(h\ge0\).
If \(z_{\ell,\iota}w_{\ell,\iota}=1\) at every index,
\(\HH h z\), and \(\HH k w\), then \(h+k\ge0\).
\end{theorem}
\begin{proof}
Let \(H,\alpha,e_\ell\) witness the first assertion.
The bound at \(\ell=2\) and Lemma~\ref{lem:algebraic} imply that
\(F=\Q(H)\) is a number field. Put \(d=[F:\Q]\ge1\) and
\(u=N_{F/\Q}(H)\in\Q^\times\).
Its complex conjugates are obtained by restricting automorphisms of
\(\C\), by Lemma~\ref{lem:extend}. Multiplying their archimedean bounds gives
\(\log\abs u\le d\alpha\). For every prime \(\ell\), the product formula
for the field norm over the \(d\) embeddings \(F\to\overline{\Q}_\ell\),
and the same extension lemma, give
\(\log\abs u_\ell\le d e_\ell\).

Here is the rational product formula explicitly. Write \(u=A/B\) in lowest
terms with \(A\in\Z\setminus\{0\}\), \(B\in\N_{>0}\).
Unique factorization gives
\[
 \log\abs u=\sum_\ell(v_\ell(A)-v_\ell(B))\log\ell,\qquad
 \log\abs u_\ell=(v_\ell(B)-v_\ell(A))\log\ell .
\]
Only primes dividing \(AB\) contribute, so for every sufficiently large \(T\)
their sum is zero. It follows that
\[
 0=\log\abs u+\sum_{\ell<T}\log\abs u_\ell
       \le d\left(\alpha+\sum_{\ell<T}e_\ell\right)\le dh .
\]
Division by \(d>0\) proves \(h\ge0\).
For the second assertion, the height calculus gives \(\HH{h+k}1\);
apply the first.
\end{proof}
\src{New/Heights/Heights.lean}{HeightLe.nonneg}

\section{One Gram determinant and completion to a basis}\label{sec:gram}
\subsection{The two embeddings in an integral Gram matrix}
For \(j:\C\to\C_\ell\) write \(\bar j=j\circ\mathrm{conj}\).
An arbitrary automorphism of \(\C\) need not commute with complex conjugation.
The two embeddings \(j,\bar j\) must therefore be treated separately.

Let \(W\) be a \(d\)-dimensional Hermitian space, \(w=(w_1,\ldots,w_d)\)
a basis, and \(T_\vartheta:W\to\mathcal H_\vartheta\) linear maps into
Hermitian spaces. Choose a base index \(0\), assume \(T_0\) injective,
and put \(z(\vartheta)=\det(T_\vartheta^*T_\vartheta)\).
Let \(\vartheta=\theta_{\ell,\iota}\) be a chosen index for each
\((\ell,\iota)\in\mathcal I\). For each such index and each embedding \(j\),
let \(\mathcal L_\vartheta(j)\subseteq\mathcal H_\vartheta\) satisfy
\begin{equation}\label{eq:pair-integral}
 \abs{j(\ip{F}{F'})}_\ell\le1
 \quad(F\in\mathcal L_\vartheta(\bar j),\
       F'\in\mathcal L_\vartheta(j)).
\end{equation}

\begin{proposition}[one-determinant bound]\label{prop:one-det}
Suppose \(c_{\ell,\iota}\ne0\), \(\nu\in\N\),
\(\HH h{(\abs{c_{\ell,\iota}}^{-2})}\), and a real sequence \(e_\ell\)
has every finite sum at most \(\delta\).
Suppose, for every \((\ell,\iota)\) and \(j:\C\to\C_\ell\), there is
a complex \(d\)-by-\(d\) matrix \(U_j\) with
\[
 \abs{j(\det U_j)}_\ell
       \ge e^{-e_\ell}\abs{j(c_{\ell,\iota})}_\ell^\nu,\qquad
 \sum_k(U_j)_{ak}T_\vartheta w_k\in\mathcal L_\vartheta(j).
\]
If every conjugate of
\(\Gamma_0=\det(\ip{T_0w_a}{T_0w_b})_{a,b}\)
has absolute value at most \(e^\alpha\), then
\[
 \HH{\alpha+2\delta+\nu h}
       {\left(z(\theta_{\ell,\iota})/z(0)\right)} .
\]
\end{proposition}
\begin{proof}
Write \(\Gamma_\vartheta=\det(\ip{T_\vartheta w_a}{T_\vartheta w_b})\).
The identity
\(\Gamma_\vartheta=z(\vartheta)\det(\ip{w_a}{w_b})\)
follows by writing \(w\) in an orthonormal basis and taking determinants
of \(T_\vartheta^*T_\vartheta\).
In particular \(\Gamma_0\ne0\) and the desired ratio is
\(\Gamma_\vartheta/\Gamma_0\).

Fix \((\ell,\iota)\), and take \(j\) to be its map into \(\C_\ell\).
Use \(U_j\) for the second slot of the pairing and \(U_{\bar j}\) for the first.
The mixed Gram matrix is
\[
 \overline{U_{\bar j}}\,
   (\ip{T_\vartheta w_a}{T_\vartheta w_b})\,
      U_j^{\mathsf t}.
\]
Its entries, and therefore its determinant, are integral at \(j\)
by \eqref{eq:pair-integral}. Taking determinants and using the two lower
bounds for the change-of-basis matrices gives
\[
 \abs{j(\Gamma_\vartheta
                  \abs{c_{\ell,\iota}}^{2\nu})}_\ell\le e^{2e_\ell}.
\]
Here \(c\bar c=\abs c^2\); no commutation of \(j\) with conjugation was used.

Thus \(H=\Gamma_0\), archimedean exponent \(\alpha\), and prime losses
\(2e_\ell\) witness a bound \(\alpha+2\delta\) for
\((\Gamma_\vartheta/\Gamma_0)\abs c^{2\nu}\).
Multiply by the \(\nu\)-th power of the family \(\abs c^{-2}\), and
apply Lemma~\ref{lem:height-calculus}.
\end{proof}
\src{New/Heights/OneDeterminant.lean}{one_determinant}

\subsection{The cost of adding the missing directions}
Here is the precise meaning of conjugate norm estimates in the next proposition.
Realize the target of \(T_0\) in a finite coordinate space with inner product
\(\bar x^{\mathsf t}Qy\), where \(Q\) is positive definite with rational entries.
Put \(\tilde\sigma=\mathrm{conj}\circ\sigma\circ\mathrm{conj}\).
For a list \(v\) let \(C_\sigma(v)\) be the Gram determinant of its
coordinatewise image under \(\sigma\). Then
\begin{equation}\label{eq:mixed-gram}
 \abs{\sigma(\det\operatorname{Gram}(v))}
 \le\sqrt{C_{\tilde\sigma}(v)\,C_\sigma(v)}.
\end{equation}
Indeed the matrix on the left is the mixed Gram matrix of
\(\tilde\sigma(v)\) and \(\sigma(v)\), since
\(\overline{\tilde\sigma(x)}=\sigma(\bar x)\) and \(Q\) is rational.
Its determinant is the inner product of the exterior products of these
two lists, so the ordinary Cauchy--Schwarz inequality proves the bound.
That exterior-power identity follows by expanding both sides as the
alternating sum over permutations. This also proves it for dependent lists.

In this coordinate setting \(\vnorm v_\sigma\) means
\(\vnorm{\sigma(v)}\), and squared covolume means \(C_\sigma\).
The same argument applies to any specified coordinate model satisfying
\eqref{eq:mixed-gram}; using it for the source's automorphic forms
requires their stated class-sum and Weyl metrics.

\begin{proposition}\label{prop:completion}
Let \(y_1,\ldots,y_n\in W\) be independent, and let a set \(S\subseteq W\)
span \(W\). Suppose the squared covolume of \(T_0y\), at every conjugate
metric under consideration, is at most \(e^V\), and
\(\vnorm{T_0s}_\sigma^2\le e^A\) for \(s\in S\) at each of those metrics.
Complete \(y\) to a basis with \(\nu=d-n\) elements \(s_1,\ldots,s_\nu\) of \(S\).
Suppose the first \(n\) vectors have integral combinations of determinant at least
\(e^{-e_\ell}\) at each \(j\), and the added vectors become integral after
multiplication by \(c_{\ell,\iota}\).
Under \eqref{eq:pair-integral}, if every finite sum of \(e_\ell\) is at most
\(\delta\) and \(\HH h{\abs c^{-2}}\), then
\[
 \HH{V+2\delta+(d-n)(A+h)}
       {\left(z(\theta_{\ell,\iota})/z(0)\right)} .
\]
\end{proposition}
\begin{proof}
The images of \(S\) span \(W/\Span(y_1,\ldots,y_n)\).
Choose a basis of the quotient from them and lift it to the indicated \(s_i\).
This proves existence of the completion and \(\nu=d-n\).

For any positive definite Hermitian form, the Gram determinant of the
completed list is at most the Gram determinant of the initial list times
the product of the squared norms of the added vectors.
Indeed Gram--Schmidt identifies the determinant with the product of the
squared orthogonal residuals; projecting onto the orthogonal complement
of the earlier span never increases a norm.
Apply this separately at \(\sigma\) and \(\tilde\sigma\), and then use
\eqref{eq:mixed-gram}. Every conjugate of the completed base Gram determinant
has absolute value at most \(e^{V+\nu A}\), as required by
Proposition~\ref{prop:one-det}.

Adjoin \(c_{\ell,\iota}I_\nu\) as a second block to the matrix for the
first \(n\) vectors. Its determinant has local absolute value at least
\(e^{-e_\ell}\abs{j(c_{\ell,\iota})}_\ell^\nu\).
Apply Proposition~\ref{prop:one-det} with \(\alpha=V+\nu A\).
\end{proof}

In the application the family \(S\) comes from the source's crude integral
rows followed by projection onto the good row space.
Its matrices for the initial \(y\) can depend on \(\ell,j,\vartheta\).
The sequence of losses \(e_\ell\) and its bound \(\delta\) cannot.
This is the quantifier order of the source's \(\operatorname{IndexLe}\).


\section{The comparison operator for one lift}\label{sec:comparison}
We first prove the linear algebra. It explains which geometric data the
construction of comparison forms must supply.

Let \(u:U\to V\) be an injective linear map of finite-dimensional Hermitian
spaces, \(\dim U=d\), let \(g:T\to V\) be another map, put \(L=\im u\),
and let \(\kappa>0\). Define
\[
 C=\kappa P_Lgg^*|_L,\qquad
 \Delta=\det C,\qquad \lambda=\det(u^*u).
\]
\begin{lemma}[compressed Gram identity]\label{lem:compressed}
If \(A:T\to U\) is surjective and \(u^*g=\gamma A\) with \(\gamma\ne0\),
then
\begin{equation}\label{eq:compressed}
 \lambda\Delta
       =\kappa^d\abs\gamma^{2d}\det(AA^*)\ne0.
\end{equation}
In particular \(g^*|_L\) is injective.
\end{lemma}
\begin{proof}
Taking adjoints gives \(g^*u=\bar\gamma A^*\).
Write \(w=g^*|_L\). Its adjoint is \(P_Lg\), so \(C=\kappa w^*w\).
Regarding \(u:U\to L\) as a bijection and taking determinants in
\((wu)^*(wu)=u^*w^*wu\), we obtain
\[
 \det((g^*u)^*(g^*u))
       =\det(u^*u)\det(w^*w).
\]
Multiplication by \(\kappa^d\) and substitution
\(g^*u=\bar\gamma A^*\) give \eqref{eq:compressed}.
Surjectivity of \(A\) makes \(A^*\) injective; thus \(AA^*\) is
positive definite. All factors on the right are nonzero.
If \(g^*z=0\) for \(z\in L\), then \(Cz=0\); invertibility of \(C\)
gives \(z=0\).
\end{proof}
\src{New/Rows/LiftVolumeHeight.lean}{adjoint_comp_eq_conj_smul; det_compressedGramOp}

\subsection{The local determinant}
\begin{lemma}[a finite family of local coordinate functionals]\label{lem:auerbach}
Let \(L\) be a \(d\)-dimensional complex vector space with transported
absolute value from \(\iota:\C\simeq\overline{\Q}_\ell\).
Let finitely many linear functionals \(l_a\) separate its points.
One can choose \(d\) of them as the coordinates of a basis \(z_1,\ldots,z_d\)
such that \(\abs{\iota(l_a(z_b))}_\ell\le1\) for every \(a,b\).
\end{lemma}
\begin{proof}
Write the functionals as rows in any basis. The separation assumption
says that the row rank is \(d\). Choose a nonzero \(d\)-row determinant
of greatest transported absolute value, among the finitely many choices.
Let \(z_b\) be the basis dual to these rows. By Cramer's rule, each
\(l_a(z_b)\) is a determinant obtained by replacing one chosen row,
divided by the chosen determinant. The numerator is either zero or
one of the same minors, up to sign, so the quotient has absolute value
at most \(1\). For \(d=0\) use the empty basis.
\end{proof}

\begin{proposition}[local comparison determinant]\label{prop:local-compressed}
Assume \(g^*|_L\) injective and fix finitely many \(f_a\in T\) such that
\(l_a(z)=\ip{f_a}{g^*z}\) separate \(L\).
Suppose each \(g f_a\) is integral at every local embedding, and whenever
\(\abs{\iota(l_a(z))}_\ell\le1\) for all \(a\), the form
\(g g^*(\kappa z)\) is integral at \(\iota\) and tested against \(L^\perp\).
If \(\operatorname{Den}_\iota(O,e_0,L;\beta)\) holds, then
\[
 \abs{\iota(\Delta)}_\ell\le\exp(d\,r_O\,\beta).
\]
\end{proposition}
\begin{proof}
Take the basis from Lemma~\ref{lem:auerbach}, with its selected coordinate
functionals \(l_{a_1},\ldots,l_{a_d}\). The matrix of \(C\) in this basis
has entries
\[
 l_{a_i}(Cz_b)
 =\ip{g f_{a_i}}{P_Lg g^*(\kappa z_b)}.
\]
This follows by moving \(g^*\) across the inner product and using
linearity in its second slot. Its first forms are integral everywhere;
the second forms are integral and tested by the assumption and the
bounds on all the \(l_a(z_b)\).
Apply \eqref{eq:den}. The determinant of a matrix of \(C\) in any basis
is \(\Delta\).
\end{proof}
\src{New/Rows/LiftVolumeHeight.lean}{norm_det_compressedGramOp_padic_le}

For the source's comparison forms the \(f_a\) are explicit Fock-dual
coordinate tables, moved to integral representatives of the finite
classes. Their functionals are the hook number times the integral
coordinate evaluations. The preceding proposition supplies the
determinant estimate once that coordinate construction has been proved.
It does not replace that construction with an arbitrary orthonormal
basis: such a basis need not be integral.

\subsection{The archimedean determinant at every automorphism}
The following version keeps track of noncontinuous automorphisms.
Suppose \(T\) has a rational positive definite coordinate pairing.
For each \(\sigma\in\Aut(\C)\), suppose \(V\) has a specified norm
\(\vnorm{\cdot}_\sigma\) with
\[
 \abs{\sigma(\ip{x}{y})}
        \le\vnorm x_{\tilde\sigma}\vnorm y_\sigma,
 \qquad \tilde\sigma=\mathrm{conj}\circ\sigma\circ\mathrm{conj}.
\]
Assume \(\kappa\) is positive rational and
\[
 \kappa\vnorm{g\zeta}_\sigma^2
       \le e^A\vnorm{\sigma(\zeta)}_T^2
 \quad(\zeta\in T,\ \sigma\in\Aut(\C)).
\]
Finally suppose
\(\vnorm{P_Lg g^*z}_\sigma\le\vnorm{g g^*z}_\sigma\) for \(z\in L\).

\begin{proposition}\label{prop:arch-compressed}
Under these assumptions and injectivity of \(g^*|_L\),
\(\abs{\sigma(\Delta)}\le e^{dA}\) for every \(\sigma\in\Aut(\C)\).
\end{proposition}
\begin{proof}
It suffices to bound each conjugated eigenvalue of \(C\), since its
determinant is their product with multiplicity.
Let \(Cz=\alpha z\) with \(z\ne0\), and put \(\eta=g^*z\ne0\).
In the coordinates of \(T\) set
\[
 \eta^\dagger
    =\mathrm{conj}\circ\sigma^{-1}\circ\mathrm{conj}\circ\sigma(\eta).
\]
Then \(\tilde\sigma(\eta^\dagger)=\sigma(\eta)\), and, since the coordinate
pairing is rational,
\(\sigma(\ip{\eta^\dagger}{\eta})=\vnorm{\sigma(\eta)}_T^2=:S>0\).
Pair \(Cz=\alpha z\) with \(g\eta^\dagger\).
Adjointness gives
\[
 \alpha\,\ip{\eta^\dagger}{\eta}
       =\kappa\,\ip{g\eta^\dagger}{P_Lg\eta}.
\]
After applying \(\sigma\) and taking absolute values, the right side is at most
\[
 \kappa\,\vnorm{g\eta^\dagger}_{\tilde\sigma}
              \vnorm{P_Lg\eta}_\sigma
 \le \kappa
       \sqrt{\frac{e^A S}{\kappa}}
       \sqrt{\frac{e^A S}{\kappa}}
 =e^A S .
\]
Divide by \(S>0\) to get \(\abs{\sigma(\alpha)}\le e^A\).
The determinant bound follows by multiplying over its \(d\) eigenvalues.
\end{proof}
\src{New/Rows/LiftVolumeHeight.lean}{norm_det_compressedGramOp_arch_le}

The contraction hypothesis is a geometric assertion about the tested
comparison forms and the Hecke kernels. Ordinary orthogonal projection
at the identity automorphism would not imply it at arbitrary \(\sigma\).
This is why the tested condition persists into the archimedean argument.

\subsection{The single-lift height bound}
Fix a lift index \(j\). For each twist \(\vartheta\), let
\(u_\vartheta:U_j\to V_{\vartheta,j}\) be its injective good adjoint lift,
and write \(\lambda_j(\vartheta)=\det(u_\vartheta^*u_\vartheta)\).
The source proves \(\dim U_j=d\), independent of \(j\).

\begin{sourceinput}[comparison forms and Hecke values]\label{in:comparison}
For every \(j\) there are a common table space \(T\), a surjection \(A:T\to U_j\),
maps \(g_\vartheta:T\to V_{\vartheta,j}\), nonzero scalars \(\gamma_\vartheta\),
a positive rational hook number \(\kappa_j\), and nonzero Hecke-value products
\(H_j(\vartheta)\), with
\[
 u_\vartheta^*g_\vartheta=\gamma_\vartheta A,\qquad
 \abs{\gamma_\vartheta/\gamma_1}^2=H_j(\vartheta)/H_j(1).
\]
They satisfy the integral coordinate and tested hypotheses of
Proposition~\ref{prop:local-compressed} for every permitted local embedding,
and the conjugate norm and contraction hypotheses of
Proposition~\ref{prop:arch-compressed}, at \(\vartheta=1\), with \(A=2rB_0\)
there. At a fit family of twists one also has
\[
 \HH{10m x^{999}+k(6B_0-\mathcal U)}
       {\left(H_j(\theta_{\ell,\iota}^j)/H_j(1)\right)^{-1}},
\]
where \(\mathcal U\) is the source's cost of one direction.
\end{sourceinput}
\src{New/Rows/ComparisonForms.lean}{exists_comparison_forms}
\src{New/Rows/HeckeValueHeight.lean}{heightLe_inv_heckeValueProd_div}

\begin{proposition}[height of one lift-volume ratio]\label{prop:one-lift}
Assume Input~\ref{in:comparison} and a fit family \(\theta\). If
\[
 3rB_0+10m x^{999}+6kB_0\le k(\log b)B_0,
\]
then
\[
 \HH{d k((\log b)B_0-\mathcal U)}
       {\left(\lambda_j(1)/\lambda_j(\theta_{\ell,\iota}^j)\right)}.
\]
\end{proposition}
\begin{proof}
Let \(\Delta(\vartheta)\) be the compressed determinant above.
Lemma~\ref{lem:compressed} shows it is nonzero and gives
\[
 \frac{\lambda_j(1)}{\lambda_j(\vartheta)}
  =\frac{\Delta(\vartheta)}{\Delta(1)}
       \left(\frac{H_j(\vartheta)}{H_j(1)}\right)^{-d}.
\]
Indeed the common factor \(\kappa_j^d\det(AA^*)\) cancels and the scalar
ratio is \(\abs{\gamma_1/\gamma_\vartheta}^{2d}\).

For the first ratio, use witness \(H=\Delta(1)\) in
Definition~\ref{def:height}. Proposition~\ref{prop:arch-compressed} gives
archimedean exponent \(2drB_0\). The local proposition and fitness give
prime losses \(dr\gamma_\ell\), whose finite sums are at most \(drB_0\).
Its height bound is therefore \(3drB_0\).
The other factor has bound \(d(10m x^{999}+k(6B_0-\mathcal U))\)
by Input~\ref{in:comparison} and the power rule.
Add, and use the displayed numerical hypothesis.
\end{proof}
\src{New/Rows/LiftVolumeHeight.lean}{heightLe_liftVolume_ratio}

\section{Theorem C: the complete final deduction}\label{sec:C}
We now give the final deduction with all signs and powers explicit.
Its source-specific inputs are stated separately so that completeness
of the deduction is not confused with completeness of their translation.

For the fixed data \(D\), put
\[
 d=\dim W,\qquad
 M=mk\mathcal U,\qquad L=(\log b)\mathcal B.
\]
Here \(W\) is the good row space.
For an \(m\)-fold twist \(\vartheta\), let
\(\Theta_\vartheta:W\to\mathcal H_\vartheta\) be the chain of theta lifts
at the prescribed level, and set
\(z(\vartheta)=\det(\Theta_\vartheta^*\Theta_\vartheta)\).
The quantities \(z(\vartheta)\) are already squared volumes.
Their squares appearing below are fourth powers of volume factors;
one must not remove the extra factor of \(2\) in the height calculation.

\begin{sourceinput}[factorization of the chain]\label{in:factorization}
All the maps \(\Theta_\vartheta\) are injective. There are nonzero complex
numbers \(c(\vartheta)\) such that
\[
 z(\vartheta)=\abs{c(\vartheta)}^d
                       \prod_{j=1}^m\lambda_j(\vartheta^j).
\]
At a fit family, the correction family has bound
\[
 \HH{2L}{\left(\abs{c(\theta_{\ell,\iota})/c(1)}^{-2}\right)}.
\]
\end{sourceinput}
\src{New/Rows/RowVolume.lean}{Data.volume_iota_eq_prod_liftVolume}

The determinant part of the factorization is elementary once its geometric
pairing identities are proved. If a common Hecke element \(T_{\vartheta,i}\)
acts on both spaces and the single-lift pairing is
\(A_{i+1}\delta(\vartheta)\ip{x}{T_{\vartheta,i}y}\), its determinant is
\((A_{i+1}\delta(\vartheta))^d\det T_{\vartheta,i}\).
The chain pairing gives a scalar to the power \(d\) times
\(\prod_i\det T_{\vartheta,i}\) on the row space.
Transfer identifies those determinants across the spaces.
Dividing the two products gives the displayed formula.
The source must also control the scalar \(c\), using the chain
normalization and the norm of the modulus. Those pairing identities,
transfer, and scalar estimates remain part of this input's translation.

\begin{sourceinput}[crude completion rows]\label{in:crude}
There is a spanning set \(S\) of the row space with base norms at most
\(\exp(2x_1^2b^2)\) in squared conjugate norm.
For each fit family, its twisted images become integral, at the primes
prime to the twist's conductor, after multiplication by a common
nonzero scalar \(\gamma_{\ell,\iota}\) with
\[
 \HH{4m^2t x^{999}}{\left(\abs{\gamma_{\ell,\iota}}^{-2}\right)}.
\]
Their integral pairings satisfy \eqref{eq:pair-integral}.
The metrics satisfy the mixed determinant comparison of
Section~\ref{sec:gram}.
\end{sourceinput}
\src{New/Exchange/CrudeRows.lean}{Data.crude_rows}
\src{New/Endgame/LongForms.lean}{Data.heightLe_volume_rows}

For an independent family \(y_1,\ldots,y_n\), its index bound
\(\operatorname{Index}(y)\le\delta\) means that a single real sequence
\(e_\ell\) has every finite sum at most \(\delta\), and for every prime,
every local embedding \(j\), and every twist \(\vartheta\) prime to that
prime, there is a matrix \(U_{\ell,j,\vartheta}\), integral at \(j\), with
\[
 \abs{j(\det U_{\ell,j,\vartheta})}_\ell\ge e^{-e_\ell},
 \qquad
 \sum_k(U_{\ell,j,\vartheta})_{ak}\Theta_\vartheta y_k
                   \text{ integral at }j.
\]
The sequence is chosen before the prime, embedding, and twist.
The matrix is chosen after them.
\src{New/Endgame/ShortAndEnough.lean}{Data.IndexLe}

The numerical hypotheses in the next statement hold for admissible data
by Proposition~\ref{prop:C-scales}.

\begin{theorem}[Theorem C, from the displayed inputs]\label{thm:C}
Assume that a fit family exists, Inputs~\ref{in:comparison},
\ref{in:factorization}, and \ref{in:crude}, and the scale comparisons
\[
 b\le L,\qquad
 2x_1^2b^2+4m^2t x^{999}\le b^3/4,\qquad
 3rB_0+10m x^{999}+6kB_0\le k(\log b)B_0 .
\]
Let \(y_1,\ldots,y_n\) be independent rows with
\[
 n\ge(1-4/b^2)d,\qquad \operatorname{Index}(y)\le\delta,
\]
and suppose that the squared covolume of \(\Theta_1y\) is at most
\(e^V\) at every conjugate metric. Then
\begin{equation}\label{eq:C}
 d(M-3L)\le V+2\delta .
\end{equation}
\end{theorem}
\begin{proof}
Choose a fit family \(\theta\).
By Proposition~\ref{prop:one-lift}, the product of its \(m\)
single-lift ratios has height at most
\[
 \sum_{j=1}^m d k((\log b)B_0-\mathcal U)
       =d(L-M),
\]
since \(mkB_0=\mathcal B\) and \(mk\mathcal U=M\).
Square that family and multiply by the \(d\)-th power of the correction
family in Input~\ref{in:factorization}.
The volume factorization and the height calculus yield
\begin{equation}\label{eq:height-lower}
 \HH{2d(2L-M)}
       {\left(z(1)^2/z(\theta_{\ell,\iota})^2\right)}.
\end{equation}
More explicitly the single-lift product contributes \(2d(L-M)\),
and the correction contributes \(2dL\). All denominators are nonzero
by injectivity and the factorization.

For the other direction use Proposition~\ref{prop:completion},
Input~\ref{in:crude}, and the index matrices for \(y\). The number of
added rows is \(\nu=d-n\le4d/b^2\). Their total cost is at most
\[
 \nu(2x_1^2b^2+4m^2t x^{999})\le db.
\]
Consequently
\begin{equation}\label{eq:height-upper}
 \HH{V+2\delta+db}
       {\left(z(\theta_{\ell,\iota})/z(1)\right)}.
\end{equation}
Both estimates use the same fixed admissible data and the same family
of twists, with the same integral structures and metrics.

The square of the family in \eqref{eq:height-upper} times that in
\eqref{eq:height-lower} is the constant family \(1\).
By Theorem~\ref{thm:height-nonneg},
\[
 0\le 2(V+2\delta+db)+2d(2L-M).
\]
Rearrangement gives
\(d(M-2L)\le V+2\delta+db\).
Subtract \(dL\) and use \(b\le L\):
\[
 d(M-3L)\le V+2\delta+d(b-L)\le V+2\delta .
\]
No sign assumption on \(M\) was used.
\end{proof}
\src{New/Endgame/LongForms.lean}{Data.integral_forms_long}

In the Lean source the inputs and the numerical scale comparisons are
theorems, so its statement of C has no additional hypotheses of this kind.
In this prose revision they remain visible until their proofs and object
definitions have been supplied. The source's completeness and this
revision's partial coverage are different facts.

\section{What still has to be written}\label{sec:remaining}
The polynomial estimates and finite-dimensional deductions above are full
proofs under their displayed hypotheses. They close several intermediate
exposition obligations, but they do not close the source-specific
arithmetic and geometry.

Sections~\ref{sec:global-selmer}--\ref{sec:pencil-twist} prove the
global Selmer construction under explicit balanced-row and value-family
hypotheses. Input~\ref{in:twist} still requires their local verification
for the automorphic rows and the Hecke values. The odd Tate-class estimate must be written with its
precise local conditions and its treatment of small primes.
The denominator section must then construct the triangular auxiliary
prime system, the trace ideal, the Ribet element with its explicit
valuation, and the integral patching operator of
Input~\ref{in:denominator}.

The geometric section must construct the comparison forms, their
Fock-dual tables and integral gauges, the Hecke-value height estimate,
the chain scalar, and the crude integral rows. It must prove the
archimedean estimates at the same data, including the tested projection
contraction. Section~\ref{sec:numerics} has discharged the budget and
Theorem C scale comparisons from admissibility. The construction of
admissible data and the further estimates needed for A and B remain.

After these chapters, A and B, the weighted prime inequality, the
good-window ending, and the five final analytic statements must be
integrated and checked against this edition.
Earlier manuscript chapters are material to compare, not evidence that
these compact-edition obligations are already discharged.
There is no defensible numerical percentage of the difficult mathematics
translated merely from the number of pages in either edition.

\begin{thebibliography}{9}
\bibitem{compact}
L. Alp\"oge and Claude; compact rewrite by L. Sernau and collaborators,
\emph{Compact formal development}, commit
\texttt{72fead102f663d5c7a407964a13dc0cf9d977207},
\href{https://github.com/lsernau/grh-lean/tree/72fead102f663d5c7a407964a13dc0cf9d977207/grh-compact}{grh-compact source}.
\bibitem{readback}
B. Naskr\k{e}cki, \emph{The condensed Lean proof of RH and GRH: a reading},
working draft, commit \texttt{82252d4e71d9f4db30a38e370dba78283ab3a921},
\href{https://github.com/nasqret/riemann-rev/tree/82252d4e71d9f4db30a38e370dba78283ab3a921/writeup}{source and individual notes}.
The notes were prepared with AI readers and writers and are not a
human-refereed exposition.
\bibitem{smithGW}
A. Smith, \emph{Galois cohomology and the Greenberg--Wiles formula:
CRM minicourse}, 6 May 2026, Theorem 1.8,
\href{https://www.asmith-math.org/CRMnotes.pdf}{author's notes}.
\bibitem{milneADT}
J. S. Milne, \emph{Arithmetic Duality Theorems}, second edition,
BookSurge, 2006, Part I,
\href{https://www.jmilne.org/math/Books/ADT.pdf}{author's edition}.
\bibitem{principia}
\emph{Unpacking the proof of RH}, Principia project,
\href{https://principia-math.com/projects/d16ab288ed353bc5b10e292ace0ddbc1/paper}{paper record},
and compact-edition compiled snapshot \texttt{sc72fead10}.
\end{thebibliography}
\end{document}
